\chapter{Chapter06 — Junta Learning Reduction}
%%% generated by the blueprint pipeline from FABL.Chapter06.LearningAndTesting.JuntaLearningReduction
\section{Overview}

Declarations of the Lean module \texttt{FABL.Chapter06.LearningAndTesting.JuntaLearningReduction}.

\section{Declarations}

\begin{definition}[Constant or nonconstant decision]
\lean{FABL.JuntaConstantDecision}
\statementsource{Analysis-of-Boolean-Functions-by-Ryan-ODonnell}{pdf-ffbe9c7b3306-p168-b001}
\label{FABL.JuntaConstantDecision}
\leanok
\uses{FABL.Sign}
A constant-or-nonconstant decision is either a declaration that the target is the
specified sign in $\{\pm1\}$, or a declaration that the target is nonconstant.
\end{definition}

\begin{definition}[Correctness of a constant decision]
\lean{FABL.JuntaConstantDecision.IsCorrect}
\label{FABL.JuntaConstantDecision.IsCorrect}
\leanok
\uses{FABL.BooleanFunction, FABL.JuntaConstantDecision, FABL.Sign}
Let $n$ be a natural number, and let $f$ be a function from $
egthinspace	ext{?}$ to $
egthinspace	ext{?}$.
\end{definition}

\begin{definition}[Constant decision value in $\mathbb F_2$]
\lean{FABL.JuntaConstantDecision.constantValueF₂?}
\leanok
\uses{FABL.JuntaConstantDecision, FABL.binarySignEquiv, FABL.𝔽₂}
\statementsource{Analysis-of-Boolean-Functions-by-Ryan-ODonnell}{pdf-ffbe9c7b3306-p168-b001}
Associate to a constant-or-nonconstant decision an optional element of $\mathbb F_2$:
for a constant sign, return the bit corresponding to that sign under the sign-to-bit
correspondence; for a nonconstant decision, return no value.
\end{definition}

\begin{definition}[Junta constant accuracy]
\lean{FABL.juntaConstantAccuracy}
\label{FABL.juntaConstantAccuracy}
\leanok
\uses{FABL.PositiveLearningParameter}
For a natural number $k$, define the accuracy parameter to be the rational number
$1/2^{k+2}$, regarded as an element of the interval $(0,1/2]$.
\end{definition}

\begin{theorem}[Rational accuracy as a real number]
\lean{FABL.juntaConstantAccuracy_cast}
\label{FABL.juntaConstantAccuracy_cast}
\leanok
\uses{FABL.juntaConstantAccuracy}
\difficulty{2}
For every natural number $k$, the real number represented by the rational junta accuracy
$1/2^{k+2}$ is $1/2^{k+2}$.
\end{theorem}

\begin{definition}[Junta constant threshold]
\lean{FABL.juntaConstantThreshold}
\label{FABL.juntaConstantThreshold}
\leanok
\uses{FABL.juntaConstantAccuracy}
For a natural number $k$, define the threshold to be $1-1/2^{k+1}$.
\end{definition}

\begin{theorem}[Real form of the testing threshold]
\lean{FABL.juntaConstantThreshold_cast}
\label{FABL.juntaConstantThreshold_cast}
\leanok
\uses{FABL.juntaConstantThreshold}
\difficulty{1}
For every natural number $k$, the real number represented by the rational
constant-testing threshold is $1-2/2^{k+2}$.
\end{theorem}

\begin{definition}[Junta constant decision rule]
\lean{FABL.decideJuntaConstant}
\statementsource{Analysis-of-Boolean-Functions-by-Ryan-ODonnell}{pdf-ffbe9c7b3306-p168-b001}
\label{FABL.decideJuntaConstant}
\leanok
\uses{FABL.JuntaConstantDecision, FABL.juntaConstantThreshold}
Let $k$ be a natural number, let $a\in\mathbb Q$ be an estimate, and write $t_k$ for the
junta constant threshold at $k$. If $t_k<a$, declare the constant target $+1$;
otherwise, if $a<-t_k$, declare the constant target $-1$; in all other cases, declare
the target nonconstant.
\end{definition}

\begin{theorem}[Zero-junta functions are constant]
\lean{FABL.exists_eq_const_of_isKJunta_zero}
\label{FABL.exists_eq_const_of_isKJunta_zero}
\leanok
\uses{FABL.BooleanFunction, FABL.IsKJunta, FABL.Sign}
\difficulty{4}
Let $f$ be a sign-valued function on the Boolean cube that depends on at most zero
coordinates. Then there is a sign value to which $f$ is identically equal.
\end{theorem}

\begin{theorem}[Mean bound for a nonconstant junta]
\lean{FABL.abs_mean_le_one_sub_inv_two_pow_of_isKJunta_of_nonconstant}
\label{FABL.abs_mean_le_one_sub_inv_two_pow_of_isKJunta_of_nonconstant}
\leanok
\uses{FABL.BooleanFunction, FABL.BooleanFunction.toReal, FABL.IsKJunta, FABL.Sign, FABL.exists_eq_const_of_isKJunta_zero, FABL.fourierDegree, FABL.fourierDegree_toReal_le_of_isKJunta, FABL.inv_two_pow_pred_le_variance_of_fourierDegree_le, FABL.mean, FABL.variance_eq_four_mul_probabilities}
\difficulty{6}
Let $n,k\in\mathbb N$, and let $f\colon\{-1,1\}^n\to\{-1,1\}$ depend on at most $k$
coordinates and be nonconstant. For the uniform mean $\mathbb E
f=2^{-n}\sum_{x\in\{-1,1\}^n}f(x)$, one has $|\mathbb E f|\leq 1-2^{-k}$.
\end{theorem}

\begin{theorem}[Accuracy parameter and inverse power identity]
\lean{FABL.inv_two_pow_eq_four_mul_juntaConstantAccuracy}
\label{FABL.inv_two_pow_eq_four_mul_juntaConstantAccuracy}
\leanok
\uses{FABL.juntaConstantAccuracy, FABL.juntaConstantAccuracy_cast}
\difficulty{3}
For every natural number $k$, let $\eta_k$ be the rational number $2^{-(k+2)}$, viewed
as the first coordinate of the learning parameter defined from $k$. Then, as real
numbers, $$(2^{-1})^k=4\eta_k.$$
\end{theorem}

\begin{theorem}[Correctness of the junta constant test]
\lean{FABL.decideJuntaConstant_isCorrect_of_close}
\label{FABL.decideJuntaConstant_isCorrect_of_close}
\leanok
\uses{FABL.BooleanFunction, FABL.BooleanFunction.toReal, FABL.IsKJunta, FABL.JuntaConstantDecision.IsCorrect, FABL.Sign, FABL.SignCube, FABL.abs_mean_le_one_sub_inv_two_pow_of_isKJunta_of_nonconstant, FABL.decideJuntaConstant, FABL.fourierCoeff, FABL.inv_two_pow_eq_four_mul_juntaConstantAccuracy, FABL.juntaConstantAccuracy, FABL.juntaConstantAccuracy_cast, FABL.juntaConstantThreshold, FABL.juntaConstantThreshold_cast, FABL.mean, FABL.mean_eq_fourierCoeff_empty, FABL.signValue_neg_one, FABL.signValue_one}
\difficulty{6}
Let $n,k\in\mathbb{N}$, and let $f:\{−1,1\}^n\to\{−1,1\}$ depend on at most $k$
coordinates. If a rational estimate $a$ differs by less than $\eta_k=1/(2^{k+2})$ from
the mean of $f$ on the uniform sign cube, then the following test returns the correct
classification: declare $f$ identically $1$ when $1-2\eta_k<a$, declare it identically
$-1$ when $a<-(1-2\eta_k)$, and otherwise declare it nonconstant.
\end{theorem}

\begin{definition}[Junta constant testing program]
\lean{FABL.juntaConstantTestProgram}
\statementsource{Analysis-of-Boolean-Functions-by-Ryan-ODonnell}{pdf-ffbe9c7b3306-p168-b001}
\label{FABL.juntaConstantTestProgram}
\leanok
\uses{FABL.JuntaConstantDecision, FABL.LearningProgram, FABL.LearningProgram.map, FABL.PositiveLearningParameter, FABL.decideJuntaConstant, FABL.juntaConstantAccuracy, FABL.scheduledFourierCoeffEstimatorProgram}
For dimension $n$, parameter $k\in\mathbb N$, and failure parameter $\delta\in(0,1/2]$,
the program with random-example access applies the junta constant decision rule at $k$
to the output of the scheduled Fourier coefficient estimator run on the empty coordinate
set with the prescribed junta constant accuracy and failure parameter $\delta$.
\end{definition}

\begin{theorem}[Cost of the constant test program]
\lean{FABL.juntaConstantTestProgram_cost_eq}
\label{FABL.juntaConstantTestProgram_cost_eq}
\leanok
\uses{FABL.BooleanFunction, FABL.JuntaConstantDecision, FABL.LearningCost, FABL.LearningCost.toTriple, FABL.LearningCost.toTriple_injective, FABL.LearningProgram.runWithCost, FABL.LearningProgram.runWithCost_map, FABL.PositiveLearningParameter, FABL.fourierEstimatorSampleCount, FABL.juntaConstantAccuracy, FABL.juntaConstantTestProgram, FABL.scheduledFourierCoeffEstimatorProgram_cost_eq}
\difficulty{4}
Let $n,k$ be natural numbers, let the target be a function from the $n$-dimensional sign
cube to $\{\pm1\}$, and let the failure parameter lie in $(0,1/2]$. For every outcome in
the support of the constant-testing program's outcome distribution, its cost is exactly
$(m,0,2m)$, where $m$ is the Fourier estimator sample count for the accuracy specified
by $k$ and the given failure parameter.
\end{theorem}

\begin{theorem}[Failure probability for junta testing]
\lean{FABL.juntaConstantTestProgram_failureProbability_le}
\label{FABL.juntaConstantTestProgram_failureProbability_le}
\leanok
\uses{FABL.BooleanFunction, FABL.IsKJunta, FABL.JuntaConstantDecision.IsCorrect, FABL.LearningProgram.eventProbability, FABL.LearningProgram.eventProbability_mono, FABL.LearningProgram.runWithCost_map, FABL.PositiveLearningParameter, FABL.decideJuntaConstant, FABL.decideJuntaConstant_isCorrect_of_close, FABL.juntaConstantAccuracy, FABL.juntaConstantTestProgram, FABL.scheduledFourierCoeffEstimatorProgram, FABL.scheduledFourierCoeffEstimatorProgram_failureProbability_le}
\difficulty{5}
Let $n,k$ be natural numbers, let the target be a function from the $n$-dimensional sign
cube to $\{\pm1\}$, and suppose it depends on at most $k$ coordinates. For every failure
parameter $\delta\in(0,1/2]$, the probability that the constant-testing program returns
an incorrect decision about whether the target is a constant sign function is at most
$\delta$.
\end{theorem}

\begin{theorem}[Sample-count upper bound]
\lean{FABL.juntaConstantTestSampleCount_cast_le}
\label{FABL.juntaConstantTestSampleCount_cast_le}
\leanok
\uses{FABL.PositiveLearningParameter, FABL.fourierEstimatorFailureBits, FABL.fourierEstimatorSampleCount, FABL.fourierEstimatorSampleCount_cast_le, FABL.juntaConstantAccuracy}
\difficulty{4}
For every natural number $k$ and failure parameter $\delta\in(0,1/2]$, the Fourier
estimator sample count using accuracy $1/2^{k+2}$ is at most $64\cdot4^k$ times the
number of failure bits associated with $\delta$.
\end{theorem}

\begin{definition}[Indices outside a coordinate set]
\lean{FABL.JuntaFreeIndex}
\label{FABL.JuntaFreeIndex}
\leanok
Let $n$ be a natural number and let $P$ be a finite set of indices in $	ext{Fin }n$. The
indices outside $P$ are the indices $i$ in $	ext{Fin }n$ such that $i\notin P$.
\end{definition}

\begin{definition}[Assignment on a coordinate set]
\lean{FABL.JuntaFixedAssignment}
\label{FABL.JuntaFixedAssignment}
\leanok
\uses{FABL.Sign}
Let $n$ be a natural number and let $P$ be a finite set of indices in $	ext{Fin }n$. A
fixed assignment on $P$ is a function assigning to each index in $P$ a sign, that is, an
element of $\{ -1,1 \}$.
\end{definition}

\begin{definition}[Assignments on free coordinates]
\lean{FABL.JuntaFreeAssignment}
\label{FABL.JuntaFreeAssignment}
\leanok
\uses{FABL.JuntaFreeIndex, FABL.Sign}
For a nonnegative integer $n$ and a set $P$ of coordinates in $[n]=\{0,\ldots,n-1\}$, a
free-coordinate assignment is a function that assigns a sign in $\{-1,1\}$ to each
coordinate in $[n]\setminus P$.
\end{definition}

\begin{definition}[Combining fixed and free assignments]
\lean{FABL.combineJuntaAssignment}
\label{FABL.combineJuntaAssignment}
\leanok
\uses{FABL.JuntaFixedAssignment, FABL.JuntaFreeAssignment}
Let $n$ be a nonnegative integer, let $P\subseteq[n]$, and choose a sign assignment on
$P$ and a sign assignment on $[n]\setminus P$. Their combination is the sign assignment
on $[n]$ that uses the first assignment on $P$ and the second on its complement.
\end{definition}

\begin{definition}[Splitting a sign assignment]
\lean{FABL.juntaAssignmentSplitEquiv}
\label{FABL.juntaAssignmentSplitEquiv}
\leanok
\uses{FABL.JuntaFixedAssignment, FABL.JuntaFreeAssignment, FABL.Sign}
Let $n$ be a nonnegative integer and $P\subseteq[n]$. Restriction to $P$ and to
$[n]\setminus P$ gives a bijection from sign assignments on $[n]$ to pairs consisting of
a sign assignment on $P$ and a sign assignment on $[n]\setminus P$.
\end{definition}

\begin{theorem}[Splitting and recombining assignments]
\lean{FABL.juntaAssignmentSplitEquiv_symm_apply}
\label{FABL.juntaAssignmentSplitEquiv_symm_apply}
\leanok
\uses{FABL.JuntaFixedAssignment, FABL.JuntaFreeAssignment, FABL.combineJuntaAssignment, FABL.juntaAssignmentSplitEquiv}
\difficulty{4}
Let $n$ be a natural number, let $P\subseteq\{0,\ldots,n-1\}$, and assign a sign in
$\{\pm1\}$ to each index in $P$ and to each index outside $P$. Splitting these
assignments into their fixed and free parts and then recombining them gives the function
on all $n$ indices whose value at each index is the assigned sign for that index.
\end{theorem}

\begin{theorem}[Fixed coordinates of a combined assignment]
\lean{FABL.combineJuntaAssignment_apply_fixed}
\label{FABL.combineJuntaAssignment_apply_fixed}
\leanok
\uses{FABL.JuntaFixedAssignment, FABL.JuntaFreeAssignment, FABL.combineJuntaAssignment}
\difficulty{1}
Let $n$ be a natural number, let $P$ be a set of coordinates in $	ext{Fin }n$, and let
$z$ assign a sign to each coordinate in $P$ and $y$ assign a sign to each coordinate
outside $P$. The assignment formed by combining $z$ and $y$ agrees with $z$ at every
coordinate in $P$.
\end{theorem}

\begin{theorem}[Combined assignment on a free coordinate]
\lean{FABL.combineJuntaAssignment_apply_free}
\label{FABL.combineJuntaAssignment_apply_free}
\leanok
\uses{FABL.JuntaFixedAssignment, FABL.JuntaFreeAssignment, FABL.JuntaFreeIndex, FABL.combineJuntaAssignment}
\difficulty{1}
Let $n$ be a natural number, let $P$ be a set of coordinates in $in n$, and let $z$
assign a sign to each coordinate in $P$ while $y$ assigns a sign to each coordinate
outside $P$. For every coordinate $i$ outside $P$, the assignment that uses $z$ on $P$
and $y$ elsewhere takes the value $y(i)$ at $i$.
\end{theorem}

\begin{definition}[Restriction to free coordinates]
\lean{FABL.juntaRestriction}
\statementsource{Analysis-of-Boolean-Functions-by-Ryan-ODonnell}{
  pdf-ffbe9c7b3306-p075-b002,
  pdf-ffbe9c7b3306-p098-b004,
  pdf-ffbe9c7b3306-p168-b001
}
\label{FABL.juntaRestriction}
\leanok
\uses{FABL.BooleanFunction, FABL.JuntaFixedAssignment, FABL.JuntaFreeAssignment, FABL.Sign, FABL.combineJuntaAssignment}
Let $n\in\mathbb N$, let $P\subseteq\{0,\ldots,n-1\}$, let $f:\{\pm1\}^n\to\{\pm1\}$,
and fix a sign assignment $z$ on $P$. The restriction of $f$ determined by $z$ is the
function on assignments to coordinates outside $P$ that combines each such assignment
with $z$ and evaluates $f$ on the resulting full assignment.
\end{definition}

\begin{definition}[Converting fixed bits to signs]
\lean{FABL.juntaFixedAssignmentSignOfF₂}
\leanok
\uses{FABL.JuntaFixedAssignment, FABL.binarySignEquiv, FABL.𝔽₂}
Let $n$ be a nonnegative integer, let $P\subseteq[n]$, and let $z:P\to\mathbb F_2$.
Assign to each coordinate in $P$ the sign corresponding to its bit under the standard
bijection between $\mathbb F_2$ and $\{-1,1\}$; this defines a sign assignment on $P$.
\end{definition}

\begin{definition}[Converting free signs to bits]
\lean{FABL.juntaFreeAssignmentF₂OfSign}
\leanok
\uses{FABL.JuntaFreeAssignment, FABL.JuntaFreeIndex, FABL.binarySignEquiv, FABL.𝔽₂}
Let $n$ be a nonnegative integer, let $P\subseteq[n]$, and let $y$ assign a sign in
$\{-1,1\}$ to each coordinate in $[n]\setminus P$. Assign to each such coordinate its
corresponding bit under the inverse of the standard bijection between $\mathbb F_2$ and
$\{-1,1\}$, obtaining a bit assignment on $[n]\setminus P$.
\end{definition}

\begin{definition}[Combining fixed and free bit assignments]
\lean{FABL.combineJuntaF₂Assignment}
\leanok
\uses{FABL.F₂Cube, FABL.JuntaFreeIndex, FABL.𝔽₂}
Let $n$ be a nonnegative integer, let $P\subseteq[n]$, and choose a bit assignment in
$\mathbb F_2$ on $P$ and one on $[n]\setminus P$. Combining them coordinatewise gives a
bit assignment in $\mathbb F_2$ on $[n]$.
\end{definition}

\begin{theorem}[Sign encoding of combined assignments]
\lean{FABL.binaryCubeSignEquiv_combineJuntaF₂Assignment}
\leanok
\uses{FABL.JuntaFreeIndex, FABL.binaryCubeSignEquiv, FABL.binaryCubeSignEquiv_apply, FABL.binarySignEquiv, FABL.combineJuntaAssignment, FABL.combineJuntaF₂Assignment, FABL.juntaFixedAssignmentSignOfF₂, FABL.𝔽₂}
\difficulty{4}
Let $P$ be a set of coordinates in $\{0,\ldots,n-1\}$, let $z$ assign a bit to each
coordinate in $P$, and let $y$ assign a bit to each coordinate outside $P$. Encoding the
binary vector formed from these assignments coordinatewise as signs gives the same sign
vector as first encoding each assignment as signs and then combining the two sign
assignments on $P$ and its complement.
\end{theorem}

\begin{definition}[Agreement with a fixed assignment]
\lean{FABL.MatchesJuntaAssignment}
\statementsource{Analysis-of-Boolean-Functions-by-Ryan-ODonnell}{pdf-ffbe9c7b3306-p168-b001}
\label{FABL.MatchesJuntaAssignment}
\leanok
\uses{FABL.JuntaFixedAssignment}
Let $n\in\mathbb N$, let $P\subseteq\{0,\ldots,n-1\}$, let $z$ assign a sign to every
coordinate in $P$, and let $x\in\{\pm1\}^n$. The point $x$ matches $z$ on $P$ exactly
when $x_i=z_i$ for every $i\in P$.
\end{definition}

\begin{definition}[Restriction to free coordinates]
\lean{FABL.juntaFreePart}
\label{FABL.juntaFreePart}
\leanok
\uses{FABL.JuntaFreeAssignment}
Let $n$ be a nonnegative integer, let $P\subseteq[n]$, and let $x:[n]\to\{-1,1\}$ be a
sign assignment. Restricting $x$ to $[n]\setminus P$ gives its free-coordinate
assignment.
\end{definition}

\begin{definition}[Restriction to fixed coordinates]
\lean{FABL.juntaFixedPart}
\label{FABL.juntaFixedPart}
\leanok
\uses{FABL.JuntaFixedAssignment}
Let $n$ be a nonnegative integer, let $P\subseteq[n]$, and let $x:[n]\to\{-1,1\}$ be a
sign assignment. Restricting $x$ to $P$ gives its fixed-coordinate assignment.
\end{definition}

\begin{theorem}[Recovery of the fixed assignment]
\lean{FABL.juntaFixedPart_combineJuntaAssignment}
\label{FABL.juntaFixedPart_combineJuntaAssignment}
\leanok
\uses{FABL.JuntaFixedAssignment, FABL.JuntaFreeAssignment, FABL.combineJuntaAssignment, FABL.juntaFixedPart}
\difficulty{3}
Let $n$ be a natural number, let $P$ be a set of coordinates among $n$, and let $z$
and $y$ assign signs to the coordinates in $P$ and outside $P$, respectively. If these
assignments are combined into a sign vector, its restriction to $P$ is $z$.
\end{theorem}

\begin{theorem}[Recovery of the free assignment]
\lean{FABL.juntaFreePart_combineJuntaAssignment}
\label{FABL.juntaFreePart_combineJuntaAssignment}
\leanok
\uses{FABL.JuntaFixedAssignment, FABL.JuntaFreeAssignment, FABL.combineJuntaAssignment, FABL.juntaFreePart}
\difficulty{3}
Let $n$ be a natural number, let $P$ be a set of coordinates among $n$, and let $z$
and $y$ assign signs to the coordinates in $P$ and outside $P$, respectively. If these
assignments are combined into a sign vector, its restriction to the coordinates outside
$P$ is $y$.
\end{theorem}

\begin{theorem}[Reconstruction from Fixed and Free Coordinates]
\lean{FABL.combineJuntaAssignment_freePart_of_matches}
\label{FABL.combineJuntaAssignment_freePart_of_matches}
\leanok
\uses{FABL.JuntaFixedAssignment, FABL.MatchesJuntaAssignment, FABL.combineJuntaAssignment, FABL.juntaFreePart}
\difficulty{4}
Let $n$ be a natural number, let $P$ be a set of coordinates in $br^n$ indexed by
$	ext{Fin}(n)$, and let $z$ assign a sign to each coordinate in $P$. If a sign vector
$x\in\{-1,1\}^{[n]}$ agrees with $z$ on every coordinate in $P$, then combining $z$ on
$P$ with the restriction of $x$ to the coordinates outside $P$ recovers $x$.
\end{theorem}

\begin{definition}[Matched junta sample]
\lean{FABL.MatchedJuntaExample}
\label{FABL.MatchedJuntaExample}
\leanok
\uses{FABL.JuntaFixedAssignment, FABL.MatchesJuntaAssignment, FABL.Sign}
Let $n$ be a natural number, let $P$ be a subset of $\\{0,\ldots,n-1\}$, and let $z$
assign a sign in $\{\pm1\}$ to each index in $P$. A matched junta sample is a pair
consisting of a point $x\in\{\pm1\}^n$ and a sign, together with the condition that $x$
matches the assignment $z$ on $P$.
\end{definition}

\begin{definition}[Restriction of a matched example]
\lean{FABL.matchedJuntaRestrictionExample}
\statementsource{Analysis-of-Boolean-Functions-by-Ryan-ODonnell}{pdf-ffbe9c7b3306-p168-b001}
\label{FABL.matchedJuntaRestrictionExample}
\leanok
\uses{FABL.JuntaFixedAssignment, FABL.JuntaFreeAssignment, FABL.MatchedJuntaExample, FABL.Sign, FABL.juntaFreePart}
Let $n$ be a natural number, let $P$ be a finite set of coordinates in
$[n]=\{0,\ldots,n-1\}$, let $z$ assign a sign in $\{-1,1\}$ to each coordinate in $P$,
and let a matched example consist of a cube point $x\in\{-1,1\}^n$ and a sign, with $x$
agreeing with $z$ on $P$. Its restriction is the pair consisting of the restriction of
$x$ to coordinates outside $P$ and its sign.
\end{definition}

\begin{theorem}[Label of a restricted junta example]
\lean{FABL.matchedJuntaRestrictionExample_label}
\label{FABL.matchedJuntaRestrictionExample_label}
\leanok
\uses{FABL.BooleanFunction, FABL.JuntaFixedAssignment, FABL.MatchedJuntaExample, FABL.combineJuntaAssignment_freePart_of_matches, FABL.juntaRestriction, FABL.matchedJuntaRestrictionExample}
\difficulty{1}
Let $n$ be a natural number, let $f:\\{-1,1\\}^n\\to\\{\\pm1\\}$ be a Boolean function,
let $P\\subseteq\\{1,\\ldots,n\\}$ be a set of coordinates, and fix signs $z_i$ for
$i\\in P$. Suppose a sample consists of an input $x\\in\\{-1,1\\}^n$ and a sign label,
with $x_i=z_i$ for every $i\\in P$, and suppose its label equals $f(x)$. Then the label
equals the value of the restriction of $f$ to the free coordinates, evaluated at the
sample's free-coordinate assignment.
\end{theorem}

\begin{definition}[Indices matching a fixed assignment]
\lean{FABL.fixedMatchingIndices}
\label{FABL.fixedMatchingIndices}
\leanok
\uses{FABL.JuntaFixedAssignment}
Let $n$ be a natural number, let $P$ be a finite subset of $\{0,\ldots,n-1\}$, and let
$z$ assign a sign to each index in $P$. For a natural number $R$ and a family
$(f_i)_{i=0}^{R-1}$ of such assignments, the matching index set is
$\{i\in\{0,\ldots,R-1\}: f_i=z\}$.
\end{definition}

\begin{definition}[Sample indices matching a junta assignment]
\lean{FABL.juntaMatchingIndices}
\label{FABL.juntaMatchingIndices}
\leanok
\uses{FABL.JuntaFixedAssignment, FABL.Sign, FABL.fixedMatchingIndices, FABL.juntaFixedPart}
Let $n$ be a natural number, let $P$ be a finite subset of $\{0,\ldots,n-1\}$, and let
$z$ assign a sign to each index in $P$. Given $R$ samples
$(x_i,s_i)\in\{\pm1\}^n\times\{\pm1\}$, the matching index set consists of those $i<R$
for which $x_i$ agrees with $z$ on every index in $P$.
\end{definition}

\begin{definition}[Number of matching samples]
\lean{FABL.juntaMatchCount}
\label{FABL.juntaMatchCount}
\leanok
\uses{FABL.JuntaFixedAssignment, FABL.Sign, FABL.juntaMatchingIndices}
Let $n$ be a natural number, let $P$ be a finite subset of $\{0,\ldots,n-1\}$, and let
$z$ assign a sign to each index in $P$. For $R$ samples
$(x_i,s_i)\in\{\pm1\}^n\times\{\pm1\}$, the match count is the number of indices $i<R$
for which $x_i$ agrees with $z$ on every index in $P$.
\end{definition}

\begin{theorem}[Counting matching samples by indicators]
\lean{FABL.juntaMatchCount_eq_sum}
\label{FABL.juntaMatchCount_eq_sum}
\leanok
\uses{FABL.JuntaFixedAssignment, FABL.MatchesJuntaAssignment, FABL.Sign, FABL.fixedMatchingIndices, FABL.juntaMatchCount, FABL.juntaMatchingIndices}
\difficulty{4}
Let $n$ be a natural number, let $P$ be a finite set of coordinates among $n$, let $z$
be a sign assignment on $P$, and let $(x_i,s_i)_{i=1}^R$ be a finite batch of samples,
where each $x_i$ is a sign vector. The number of samples whose coordinates on $P$ agree
with $z$ equals the sum over $i=1, R$ of $1$ when $x_i$ agrees with $z$ on $P$, and $0$
otherwise.
\end{theorem}

\begin{definition}[Junta match indicator]
\lean{FABL.juntaMatchObservation}
\label{FABL.juntaMatchObservation}
\leanok
\uses{FABL.JuntaFixedAssignment, FABL.MatchesJuntaAssignment}
Let $n$ be a natural number, let $P$ be a finite subset of $\{0,\ldots,n-1\}$, let $z$
assign a sign to each index in $P$, and let $x\in\{\pm1\}^n$. The match observation is
$1$ if $x$ agrees with $z$ on every index in $P$, and is $0$ otherwise.
\end{definition}

\begin{theorem}[Bounds on the match observation]
\lean{FABL.juntaMatchObservation_mem_Icc}
\label{FABL.juntaMatchObservation_mem_Icc}
\leanok
\uses{FABL.JuntaFixedAssignment, FABL.MatchesJuntaAssignment, FABL.juntaMatchObservation}
\difficulty{1}
Let $n$ be a natural number, let $P$ be a finite set of coordinates among $n$, let $z$
be a sign assignment on $P$, and let $x$ be a sign vector. The observation that equals
$1$ if $x$ agrees with $z$ on $P$ and $0$ otherwise belongs to the interval $[-1,1]$.
\end{theorem}

\begin{theorem}[Probability of matching a fixed assignment]
\lean{FABL.expect_juntaMatchObservation}
\label{FABL.expect_juntaMatchObservation}
\leanok
\uses{FABL.JuntaFixedAssignment, FABL.JuntaFreeAssignment, FABL.MatchesJuntaAssignment, FABL.Sign, FABL.combineJuntaAssignment, FABL.juntaAssignmentSplitEquiv, FABL.juntaAssignmentSplitEquiv_symm_apply, FABL.juntaMatchObservation, FABL.uniformProbability_eq_singleton}
\difficulty{5}
Let $n$ be a nonnegative integer, let $P$ be a subset of $\\{0,\\ldots,n-1\\}$, and
assign a sign $z_i\\in\\{-1,1\\}$ to each $i\\in P$. For a uniformly random sign vector
$x\\in\\{-1,1\\}^n$, the probability that $x_i=z_i$ for every $i\\in P$ is $2^{-|P|}$.
\end{theorem}

\begin{theorem}[Empirical junta match frequency]
\lean{FABL.finiteUniformEmpiricalMean_juntaMatchObservation}
\label{FABL.finiteUniformEmpiricalMean_juntaMatchObservation}
\leanok
\uses{FABL.BooleanFunction, FABL.JuntaFixedAssignment, FABL.MatchesJuntaAssignment, FABL.finiteUniformEmpiricalMean, FABL.juntaMatchCount, FABL.juntaMatchCount_eq_sum, FABL.juntaMatchObservation}
\difficulty{3}
Let $n$ and $R$ be nonnegative integers, let $f:\{-1,1\}^n\to\{-1,1\}$, let $P$ be a
subset of $\{0,\ldots,n-1\}$, let $z:P\to\{-1,1\}$, and let
$x_0,\ldots,x_{R-1}\in\{-1,1\}^n$. The empirical mean of the indicator that an input
agrees with $z$ on every coordinate in $P$, evaluated on these inputs, equals $1/R$
times the number of inputs among them that agree with $z$ on $P$.
\end{theorem}

\begin{definition}[Initial segment embedding]
\lean{FABL.initialFinEmbedding}
\label{FABL.initialFinEmbedding}
\leanok
For natural numbers $M\le L$, the map from $\{0,\ldots,M-1\}$ to $\{0,\ldots,L-1\}$
sending each index to itself is an embedding.
\end{definition}

\begin{definition}[Embedding into matching indices]
\lean{FABL.fixedMatchingIndexEmbedding}
\label{FABL.fixedMatchingIndexEmbedding}
\leanok
\uses{FABL.JuntaFixedAssignment, FABL.fixedMatchingIndices, FABL.initialFinEmbedding}
Let $n$ be a natural number, let $P$ be a finite subset of $\{0,\ldots,n-1\}$, and let
$z$ assign a sign to each index in $P$. For natural numbers $M,R$, let $(f_i)_{i<R}$ be
assignments on $P$, and suppose at least $M$ of them equal $z$. Then there is an
embedding of $\{0,\ldots,M-1\}$ into $\{0,\ldots,R-1\}$ whose image consists of indices
with $f_i=z$.
\end{definition}

\begin{definition}[Embedding into matching sample indices]
\lean{FABL.juntaMatchingIndexEmbedding}
\label{FABL.juntaMatchingIndexEmbedding}
\leanok
\uses{FABL.JuntaFixedAssignment, FABL.Sign, FABL.fixedMatchingIndexEmbedding, FABL.juntaFixedPart, FABL.juntaMatchCount}
Let $n$ be a natural number, let $P$ be a finite subset of $\{0,\ldots,n-1\}$, and let
$z$ assign a sign to each index in $P$. Given natural numbers $M,R$ and samples
$(x_i,s_i)\in\{\pm1\}^n\times\{\pm1\}$ for $i<R$, suppose at least $M$ sample points
agree with $z$ on every index in $P$. Then there is an embedding of $\{0,\ldots,M-1\}$
into $\{0,\ldots,R-1\}$ whose image consists of indices of such matching samples.
\end{definition}

\begin{theorem}[Selected samples match the assignment]
\lean{FABL.juntaMatchingIndexEmbedding_matches}
\label{FABL.juntaMatchingIndexEmbedding_matches}
\leanok
\uses{FABL.JuntaFixedAssignment, FABL.MatchesJuntaAssignment, FABL.Sign, FABL.juntaMatchCount, FABL.juntaMatchingIndexEmbedding, FABL.juntaMatchingIndices}
\difficulty{3}
Let $n,R,M$ be natural numbers, let $P$ be a finite set of coordinates among $n$, let
$z$ be a sign assignment on $P$, and let $(x_r,s_r)_{r=1}^R$ be a batch of samples. If
$M$ is at most the number of samples whose coordinates on $P$ agree with $z$, then every
sample selected by the canonical order-preserving injection from the first $M$ indices
into those matching indices agrees with $z$ on $P$.
\end{theorem}

\begin{definition}[Selecting matching junta examples]
\lean{FABL.takeMatchingJuntaExamples}
\statementsource{Analysis-of-Boolean-Functions-by-Ryan-ODonnell}{pdf-ffbe9c7b3306-p168-b001}
\label{FABL.takeMatchingJuntaExamples}
\leanok
\uses{FABL.JuntaFixedAssignment, FABL.MatchedJuntaExample, FABL.Sign, FABL.juntaMatchCount, FABL.juntaMatchingIndexEmbedding, FABL.juntaMatchingIndexEmbedding_matches}
Let $n$, $M$, and $R$ be natural numbers, let $P$ be a finite set of coordinates in
$\\{0,\ldots,n-1\}$, and let $z$ assign a sign to each coordinate in $P$. Given a
sequence of $R$ labeled samples from $\{\pm1\}^n$, the construction returns a tuple of
$M$ samples that match $z$ on $P$ if at least $M$ samples in the sequence match;
otherwise, it returns no tuple.
\end{definition}

\begin{theorem}[When matching junta examples are unavailable]
\lean{FABL.takeMatchingJuntaExamples_eq_none_iff}
\label{FABL.takeMatchingJuntaExamples_eq_none_iff}
\leanok
\uses{FABL.JuntaFixedAssignment, FABL.Sign, FABL.juntaMatchCount, FABL.takeMatchingJuntaExamples}
\difficulty{3}
Let $P$ be a set of coordinates in the $n$-dimensional sign cube, let $z$ assign a sign
to each coordinate in $P$, let $M$ be a natural number, and let a finite list of $R$
examples be given, each consisting of an input and an output sign. The procedure that
selects $M$ examples whose inputs agree with $z$ on every coordinate in $P$ returns no
selection if and only if fewer than $M$ examples in the list agree with $z$ on all those
coordinates.
\end{theorem}

\begin{definition}[Junta restriction accuracy parameter]
\lean{FABL.juntaRestrictionMatchAccuracy}
\label{FABL.juntaRestrictionMatchAccuracy}
\leanok
\uses{FABL.PositiveLearningParameter}
For each natural number $k$, define the junta restriction match accuracy to be the
rational number $2^{-(k+1)}$, regarded as an element of $(0,1/2]$.
\end{definition}

\begin{theorem}[Real form of the matching accuracy]
\lean{FABL.juntaRestrictionMatchAccuracy_cast}
\label{FABL.juntaRestrictionMatchAccuracy_cast}
\leanok
\uses{FABL.juntaRestrictionMatchAccuracy}
\difficulty{2}
For every natural number $k$, the real cast of the rational matching accuracy
$2^{-(k+1)}$ equals $2^{-(k+1)}$.
\end{theorem}

\begin{definition}[Junta restriction sample count]
\lean{FABL.juntaRestrictionSampleCount}
\statementsource{Analysis-of-Boolean-Functions-by-Ryan-ODonnell}{pdf-ffbe9c7b3306-p168-b001}
\label{FABL.juntaRestrictionSampleCount}
\leanok
\uses{FABL.PositiveLearningParameter, FABL.fourierEstimatorSampleCount, FABL.juntaRestrictionMatchAccuracy}
For natural numbers $k$ and $M$ and a failure parameter $\delta$ with $0<\delta\leq
1/2$, define the sample count to be $2^{k+1}M$ plus the Fourier-estimator sample count
at accuracy equal to the junta-restriction match accuracy for $k$, with failure
parameter $\delta$.
\end{definition}

\begin{definition}[Junta restriction sample work]
\lean{FABL.juntaRestrictionSampleWork}
\statementsource{Analysis-of-Boolean-Functions-by-Ryan-ODonnell}{pdf-ffbe9c7b3306-p168-b001}
\label{FABL.juntaRestrictionSampleWork}
\leanok
\uses{FABL.PositiveLearningParameter, FABL.juntaRestrictionSampleCount}
Let $n$ be a natural number, let $P$ be a finite set of coordinates in $[n]$, let $k,M$
be natural numbers, and let the rational failure parameter $\delta$ satisfy
$0<\delta\leq 1/2$. Define the work bound as the junta restriction sample count for
$k,M,\delta$, multiplied by $|P|+1$, plus $M$.
\end{definition}

\begin{definition}[Junta restriction sample program]
\lean{FABL.juntaRestrictionSampleProgram}
\statementsource{Analysis-of-Boolean-Functions-by-Ryan-ODonnell}{pdf-ffbe9c7b3306-p168-b001}
\label{FABL.juntaRestrictionSampleProgram}
\leanok
\uses{FABL.JuntaFixedAssignment, FABL.LearningProgram, FABL.MatchedJuntaExample, FABL.PositiveLearningParameter, FABL.juntaRestrictionSampleCount, FABL.juntaRestrictionSampleWork, FABL.takeMatchingJuntaExamples}
For a natural number $n$, a finite coordinate set $P\subseteq[n]$, an assignment $z$ of
signs on $P$, natural numbers $k,M$, and a rational failure parameter $\delta$ with
$0<\delta\leq 1/2$, consider the random-example learning program that draws the
prescribed junta-restriction sample batch, charges the prescribed junta-restriction
work, and returns the result of selecting up to $M$ examples from the batch that match
$z$ on $P$.
\end{definition}

\begin{theorem}[Run with cost junta restriction sample program]
\lean{FABL.runWithCost_juntaRestrictionSampleProgram}
\statementsource{Analysis-of-Boolean-Functions-by-Ryan-ODonnell}{pdf-ffbe9c7b3306-p168-b001}
\label{FABL.runWithCost_juntaRestrictionSampleProgram}
\leanok
\uses{FABL.BooleanFunction, FABL.JuntaFixedAssignment, FABL.LearningProgram.runWithCost, FABL.PositiveLearningParameter, FABL.juntaRestrictionSampleCount, FABL.juntaRestrictionSampleProgram, FABL.juntaRestrictionSampleWork, FABL.takeMatchingJuntaExamples, FABL.uniformPMF}
\difficulty{3}
% dataset-pipeline: stub -- prose pending build_informal_statements.py
Statement pending regeneration.
\end{theorem}

\begin{definition}[Junta restriction failure set]
\lean{FABL.juntaRestrictionFailureSet}
\label{FABL.juntaRestrictionFailureSet}
\leanok
\uses{FABL.BooleanFunction, FABL.JuntaFixedAssignment, FABL.PositiveLearningParameter, FABL.juntaRestrictionSampleCount, FABL.takeMatchingJuntaExamples}
Let $n,k,M$ be natural numbers, let $f$ be a Boolean function on the $n$-dimensional
Boolean cube, let $P$ be a finite set of coordinates in $	ext{Fin}(n)$, let $z$ assign a
sign to each coordinate in $P$, and let $ailure	ext{-parameter}$ lie in $(0,1/2]$. The
junta restriction failure set consists of all sample sequences of length equal to the
junta restriction sample count for $k$, $M$, and the failure parameter for which
selecting $M$ examples matching $z$ from the samples labeled by $f$ returns no batch.
\end{definition}

\begin{theorem}[Junta restriction failure and empirical error]
\lean{FABL.juntaRestrictionFailureSet_subset_empiricalBad}
\label{FABL.juntaRestrictionFailureSet_subset_empiricalBad}
\leanok
\uses{FABL.BooleanFunction, FABL.JuntaFixedAssignment, FABL.PositiveLearningParameter, FABL.expect_juntaMatchObservation, FABL.finiteUniformEmpiricalMean, FABL.finiteUniformEmpiricalMean_juntaMatchObservation, FABL.fourierEstimatorSampleCount_pos, FABL.juntaMatchCount, FABL.juntaMatchObservation, FABL.juntaRestrictionFailureSet, FABL.juntaRestrictionMatchAccuracy, FABL.juntaRestrictionSampleCount, FABL.takeMatchingJuntaExamples, FABL.takeMatchingJuntaExamples_eq_none_iff}
\difficulty{5}
Let $n,k,M\in\mathbb{N}$, let $f:\{−1,1\}^n\to\{−1,1\}$ be a Boolean function, let
$P\subseteq\{1,\ldots,n\}$, and fix a sign assignment $z$ on $P$. Suppose $|P|\le k$,
and let $\delta$ satisfy $0<\delta\le\tfrac12$. For every sequence of $2^{k+1}M$ plus
the Fourier estimator sample count at accuracy $2^{-(k+1)}$ and failure parameter
$\delta$ cube inputs, if fewer than $M$ inputs agree with $z$ on all coordinates in $P$,
then the absolute difference between the empirical fraction agreeing with $z$ and its
uniform-cube probability is at least $2^{-(k+1)}$.
\end{theorem}

\begin{theorem}[Junta restriction sample program failure probability le]
\lean{FABL.juntaRestrictionSampleProgram_failureProbability_le}
\statementsource{Analysis-of-Boolean-Functions-by-Ryan-ODonnell}{pdf-ffbe9c7b3306-p168-b001}
\label{FABL.juntaRestrictionSampleProgram_failureProbability_le}
\leanok
\uses{FABL.BooleanFunction, FABL.JuntaFixedAssignment, FABL.LearningCost, FABL.LearningProgram.eventProbability, FABL.PositiveLearningParameter, FABL.finiteUniformEmpiricalMean, FABL.fourierEstimatorSampleCount, FABL.fourierEstimatorSampleCount_pos, FABL.juntaMatchObservation, FABL.juntaMatchObservation_mem_Icc, FABL.juntaRestrictionFailureSet, FABL.juntaRestrictionFailureSet_subset_empiricalBad, FABL.juntaRestrictionMatchAccuracy, FABL.juntaRestrictionSampleCount, FABL.juntaRestrictionSampleProgram, FABL.juntaRestrictionSampleWork, FABL.measure_finiteUniformEmpiricalMean_sub_expect_ge_le, FABL.positiveLearningParameter_toReal_mem_Ioc, FABL.runWithCost_juntaRestrictionSampleProgram, FABL.takeMatchingJuntaExamples, FABL.two_mul_exp_neg_fourierEstimatorSampleCount_le, FABL.uniformPMF}
\difficulty{5}
% dataset-pipeline: stub -- prose pending build_informal_statements.py
Statement pending regeneration.
\end{theorem}

\begin{theorem}[Correct labels on matched examples]
\lean{FABL.juntaRestrictionSampleProgram_success_labels}
\label{FABL.juntaRestrictionSampleProgram_success_labels}
\leanok
\uses{FABL.BooleanFunction, FABL.JuntaFixedAssignment, FABL.LearningCost, FABL.LearningProgram.runWithCost, FABL.MatchedJuntaExample, FABL.PositiveLearningParameter, FABL.juntaRestriction, FABL.juntaRestrictionSampleProgram, FABL.matchedJuntaRestrictionExample, FABL.matchedJuntaRestrictionExample_label, FABL.runWithCost_juntaRestrictionSampleProgram, FABL.takeMatchingJuntaExamples}
\difficulty{4}
Let $n,k,M\in\mathbb{N}$, let $f:\{−1,1\}^n\to\{−1,1\}$ be a Boolean function, let
$P\subseteq\{1,\ldots,n\}$, and fix an assignment $z$ on $P$. Let $\delta$ satisfy
$0<\delta\le\tfrac12$. If an outcome of the program's run has first component equal to a
batch of $M$ matched examples, then for every example in that batch its label equals the
value of the restriction of $f$ obtained by fixing $P$ according to $z$ and using that
example's remaining coordinates as input.
\end{theorem}

\begin{theorem}[Cost of the junta restriction program]
\lean{FABL.juntaRestrictionSampleProgram_cost_eq}
\label{FABL.juntaRestrictionSampleProgram_cost_eq}
\leanok
\uses{FABL.BooleanFunction, FABL.JuntaFixedAssignment, FABL.LearningCost, FABL.LearningProgram.runWithCost, FABL.MatchedJuntaExample, FABL.PositiveLearningParameter, FABL.juntaRestrictionSampleCount, FABL.juntaRestrictionSampleProgram, FABL.juntaRestrictionSampleWork, FABL.runWithCost_juntaRestrictionSampleProgram}
\difficulty{3}
Let $n,k,M\in\mathbb{N}$, let $f:\{−1,1\}^n\to\{−1,1\}$ be a Boolean function, let
$P\subseteq\{1,\ldots,n\}$, and fix an assignment $z$ on $P$. Let $\delta$ satisfy
$0<\delta\le\tfrac12$. Every outcome in the support of the program's run with target $f$
has cost $(S,0,S+W)$, where $S$ is the program's defined sample count and $W$ its
defined sample work for $P,k,M,\delta$.
\end{theorem}

\begin{theorem}[Upper bound on the sample count]
\lean{FABL.juntaRestrictionSampleCount_cast_le}
\label{FABL.juntaRestrictionSampleCount_cast_le}
\leanok
\uses{FABL.PositiveLearningParameter, FABL.fourierEstimatorFailureBits, FABL.fourierEstimatorSampleCount, FABL.fourierEstimatorSampleCount_cast_le, FABL.juntaRestrictionMatchAccuracy, FABL.juntaRestrictionSampleCount}
\difficulty{4}
Let $k,M\in\mathbb{N}$ and let $\delta$ satisfy $0<\delta\le\tfrac12$. The rational cast
of the junta restriction sample count is at most $2^{k+1}M+16\cdot4^k B(\delta)$, where
$B(\delta)$ is the number of failure bits used by the Fourier estimator with parameter
$\delta$.
\end{theorem}

\begin{definition}[Injection coordinate splitting equivalence]
\lean{FABL.injectionVectorSplitEquiv}
\label{FABL.injectionVectorSplitEquiv}
\leanok
Let $A$ be a type, let $M,R$ be natural numbers, and let $e:	ext{Fin}(M)	o	ext{Fin}(R)$
be an injection. The space of $R$-indexed $A$-valued vectors is equivalent to the
product of an $M$-indexed $A$-valued vector and an $A$-valued assignment on the
coordinates of $	ext{Fin}(R)$ outside the image of $e$.
\end{definition}

\begin{theorem}[Restriction along an injection]
\lean{FABL.injectionVectorSplitEquiv_fst}
\label{FABL.injectionVectorSplitEquiv_fst}
\leanok
\uses{FABL.injectionVectorSplitEquiv}
\difficulty{1}
Let $A$ be a type, let $M,R$ be nonnegative integers, let
$e:\{0,\ldots,M-1\}\hookrightarrow\{0,\ldots,R-1\}$ be an injection, and let $x$ assign
an element of $A$ to each of the $R$ indices. Under the canonical decomposition of such
assignments into values on the image of $e$ and on its complement, the first component
at index $i$ is $x(e(i))$ for every $i<M$.
\end{theorem}

\begin{theorem}[First marginal of a uniform distribution]
\lean{FABL.map_uniformPMF_fst}
\label{FABL.map_uniformPMF_fst}
\leanok
\uses{FABL.uniformPMF}
\difficulty{4}
Let $A$ and $B$ be finite nonempty sets. The image of the uniform probability mass
function on $A\times B$ under projection onto $A$ is the uniform probability mass
function on $A$.
\end{theorem}

\begin{theorem}[Uniform distribution under coordinate projection]
\lean{FABL.map_uniformPMF_injection_projection}
\label{FABL.map_uniformPMF_injection_projection}
\leanok
\uses{FABL.injectionVectorSplitEquiv, FABL.map_uniformPMF_equiv, FABL.map_uniformPMF_fst, FABL.uniformPMF}
\difficulty{3}
Let $A$ be a finite nonempty set, and let
$e:\{0,\ldots,M-1\}\hookrightarrow\{0,\ldots,R-1\}$ be an injection. For a uniformly
chosen function $x:\{0,\ldots,R-1\}\to A$, the function $i\mapsto x(e(i))$ is uniformly
distributed among all functions from $\{0,\ldots,M-1\}$ to $A$.
\end{theorem}

\begin{definition}[Selected free assignments for a junta]
\lean{FABL.selectedJuntaFreeBatch}
\label{FABL.selectedJuntaFreeBatch}
\leanok
\uses{FABL.JuntaFixedAssignment, FABL.JuntaFreeAssignment, FABL.fixedMatchingIndexEmbedding, FABL.fixedMatchingIndices}
Let $n,M$ be natural numbers, let $P$ be a finite set of coordinates in $	ext{Fin}(n)$,
and let $z$ assign signs to the coordinates in $P$. For a natural number $R$, suppose
$M$ is at most the number of indices $r	ext{ in }	ext{Fin}(R)$ whose fixed assignment
equals $z$, and let each such index have a free assignment on the coordinates outside
$P$. The selected batch is the $M$-indexed list of free assignments at the matching
indices chosen by the fixed-matching-index embedding.
\end{definition}

\begin{definition}[Labeled examples of a restricted function]
\lean{FABL.juntaRestrictionLabeledBatch}
\statementsource{Analysis-of-Boolean-Functions-by-Ryan-ODonnell}{pdf-ffbe9c7b3306-p168-b001}
\label{FABL.juntaRestrictionLabeledBatch}
\leanok
\uses{FABL.BooleanFunction, FABL.JuntaFixedAssignment, FABL.JuntaFreeAssignment, FABL.Sign, FABL.juntaRestriction}
Let $n$ be a natural number, let the target be a function from the $n$-dimensional sign
cube to signs, let $P$ be a finite set of coordinates in $[n]$, let $z$ assign signs on
$P$, and let $M$ be a natural number. Given an $M$-tuple of assignments on the
coordinates outside $P$, produce an $M$-tuple whose $i$th entry consists of that
free-coordinate assignment and the value of the target's restriction determined by $P$,
$z$, and that assignment.
\end{definition}

\begin{definition}[Labeled samples from split assignments]
\lean{FABL.juntaLabeledSamplesFromSplit}
\statementsource{Analysis-of-Boolean-Functions-by-Ryan-ODonnell}{pdf-ffbe9c7b3306-p168-b001}
\label{FABL.juntaLabeledSamplesFromSplit}
\leanok
\uses{FABL.BooleanFunction, FABL.JuntaFixedAssignment, FABL.JuntaFreeAssignment, FABL.Sign, FABL.combineJuntaAssignment}
Let $n$ be a natural number, let the target be a function from the $n$-dimensional sign
cube to signs, and let $P$ be a finite set of coordinates in $[n]$. Given $R$ pairs of
assignments, with each pair consisting of an assignment on $P$ and an assignment on its
complement, combine the two assignments to form a cube point and pair that point with
its target value; the result is an $R$-tuple of labeled samples.
\end{definition}

\begin{theorem}[Matched Examples from a Split Sample]
\lean{FABL.takeMatchingJuntaExamples_from_split}
\label{FABL.takeMatchingJuntaExamples_from_split}
\leanok
\uses{FABL.BooleanFunction, FABL.JuntaFixedAssignment, FABL.JuntaFreeAssignment, FABL.fixedMatchingIndexEmbedding, FABL.fixedMatchingIndices, FABL.initialFinEmbedding, FABL.juntaFixedPart, FABL.juntaFreePart_combineJuntaAssignment, FABL.juntaLabeledSamplesFromSplit, FABL.juntaMatchCount, FABL.juntaMatchingIndexEmbedding, FABL.juntaMatchingIndices, FABL.juntaRestriction, FABL.juntaRestrictionLabeledBatch, FABL.matchedJuntaRestrictionExample, FABL.selectedJuntaFreeBatch, FABL.takeMatchingJuntaExamples}
\difficulty{5}
Let $n,R,M$ be nonnegative integers, let $P$ be a subset of the $n$ coordinates, and let
$z$ assign a sign in $\\{ -1,1 \\}$ to each coordinate in $P$. Let a Boolean function
assign a sign to every sign vector, and let each of $R$ samples be formed by combining a
fixed assignment on $P$ with a free assignment off $P$, then labeling the combined
vector by the function. If at least $M$ of the fixed assignments equal $z$, then
selecting the first $M$ such samples and retaining their free assignments and labels
gives exactly the batch obtained by applying the function to those free assignments with
$z$ fixed on $P$; in particular, the option-valued selection is defined and equals this
batch.
\end{theorem}

\begin{theorem}[Distribution of a selected restriction batch]
\lean{FABL.map_uniformPMF_selectedJuntaRestrictionBatch}
\label{FABL.map_uniformPMF_selectedJuntaRestrictionBatch}
\leanok
\uses{FABL.BooleanFunction, FABL.JuntaFixedAssignment, FABL.JuntaFreeAssignment, FABL.fixedMatchingIndexEmbedding, FABL.fixedMatchingIndices, FABL.juntaRestrictionLabeledBatch, FABL.map_uniformPMF_injection_projection, FABL.selectedJuntaFreeBatch, FABL.uniformPMF}
\difficulty{3}
Let $n,M,R$ be natural numbers, let $f$ map sign vectors on $n$ coordinates to signs,
let $P$ be a set of coordinates, and fix a sign assignment $z$ on $P$. Given a batch of
$R$ assignments on $P$, suppose at least $M$ entries equal $z$. Select $M$ such entries
by the specified order-preserving selection, and use their corresponding free-coordinate
assignments to form a batch of labeled restrictions under $f$. When the original batch
of free-coordinate assignments is uniform, the resulting labeled batch has the same
distribution as a batch formed by choosing $M$ free-coordinate assignments uniformly and
labeling each by its restriction under $f$.
\end{theorem}

\begin{definition}[Projection of a matched batch]
\lean{FABL.projectJuntaRestrictionBatch}
\statementsource{Analysis-of-Boolean-Functions-by-Ryan-ODonnell}{pdf-ffbe9c7b3306-p168-b001}
\label{FABL.projectJuntaRestrictionBatch}
\leanok
\uses{FABL.JuntaFixedAssignment, FABL.JuntaFreeAssignment, FABL.MatchedJuntaExample, FABL.Sign, FABL.matchedJuntaRestrictionExample}
Let $n$ be a natural number, let $P$ be a finite set of coordinates in $[n]$, let $z$
assign signs on $P$, and let $M$ be a natural number. A missing optional batch remains
missing; when an $M$-tuple of labeled samples matching $z$ on $P$ is supplied, map each
sample to the pair formed from its free-coordinate restriction and its label, yielding
an optional $M$-tuple of such pairs.
\end{definition}

\begin{definition}[Bad junta restriction sample]
\lean{FABL.JuntaRestrictionSampleBad}
\label{FABL.JuntaRestrictionSampleBad}
\leanok
\uses{FABL.JuntaFixedAssignment, FABL.JuntaFreeAssignment, FABL.MatchedJuntaExample, FABL.Sign, FABL.projectJuntaRestrictionBatch}
Let $n,M$ be natural numbers, let $P$ be a finite set of coordinates in $	ext{Fin}(n)$,
let $z$ assign signs to the coordinates in $P$, and let $	ext{bad}$ be a predicate on
$M$-tuples of free assignments paired with signs. Given either no sampled batch or a
batch of $M$ examples that match $z$ on $P$, the sample is bad when projection of the
sampled batch to free assignments and labels returns no batch, or when it returns a
batch satisfying $	ext{bad}$.
\end{definition}

\begin{definition}[Junta input batch splitting equivalence]
\lean{FABL.juntaInputBatchSplitEquiv}
\label{FABL.juntaInputBatchSplitEquiv}
\leanok
\uses{FABL.JuntaFixedAssignment, FABL.JuntaFreeAssignment, FABL.combineJuntaAssignment, FABL.juntaFixedPart, FABL.juntaFreePart}
Let $n,R$ be natural numbers and let $P$ be a finite set of coordinates in
$	ext{Fin}(n)$. The space of $R$-indexed batches of points in the $n$-dimensional
Boolean cube is equivalent to the product of an $R$-indexed batch of sign assignments on
$P$ and an $R$-indexed batch of sign assignments on the coordinates outside $P$.
\end{definition}

\begin{theorem}[Reconstructing a batch from its split]
\lean{FABL.juntaInputBatchSplitEquiv_symm_apply}
\label{FABL.juntaInputBatchSplitEquiv_symm_apply}
\leanok
\uses{FABL.JuntaFixedAssignment, FABL.JuntaFreeAssignment, FABL.combineJuntaAssignment, FABL.juntaInputBatchSplitEquiv}
\difficulty{1}
Let $n,R$ be natural numbers and let $P$ be a finite set of coordinates among $n$.
Given a batch of assignments of signs to the coordinates in $P$ and a batch of
assignments to the coordinates outside $P$, applying the inverse of the batch-splitting
bijection returns the batch whose $r$th entry is the combination of the two assignments
at $r$.
\end{theorem}

\begin{definition}[Raw projected junta restriction batch]
\lean{FABL.rawProjectedJuntaRestrictionBatch}
\statementsource{Analysis-of-Boolean-Functions-by-Ryan-ODonnell}{pdf-ffbe9c7b3306-p168-b001}
\label{FABL.rawProjectedJuntaRestrictionBatch}
\leanok
\uses{FABL.BooleanFunction, FABL.JuntaFixedAssignment, FABL.JuntaFreeAssignment, FABL.Sign, FABL.projectJuntaRestrictionBatch, FABL.takeMatchingJuntaExamples}
Let $n$, $M$, and $R$ be natural numbers, let $f$ be a sign-valued Boolean function on
the $n$-dimensional sign cube, let $P$ be a set of coordinates in $[n]$, and fix a sign
assignment on $P$. Given a sequence of $R$ inputs from the sign cube, the raw projected
junta restriction batch is either undefined or a sequence of $M$ pairs, each consisting
of an assignment on the coordinates outside $P$ and a sign, obtained by selecting
matching labeled examples and projecting them to the restriction determined by the fixed
assignment.
\end{definition}

\begin{theorem}[Distribution of a raw projected restriction batch]
\lean{FABL.map_uniformPMF_rawProjectedJuntaRestrictionBatch}
\label{FABL.map_uniformPMF_rawProjectedJuntaRestrictionBatch}
\leanok
\uses{FABL.BooleanFunction, FABL.JuntaFixedAssignment, FABL.JuntaFreeAssignment, FABL.combineJuntaAssignment, FABL.juntaInputBatchSplitEquiv, FABL.map_uniformPMF_equiv, FABL.rawProjectedJuntaRestrictionBatch, FABL.uniformPMF, FABL.uniformPMF_bind_map_pair}
\difficulty{4}
Let $n,M,R$ be natural numbers, let $f$ map sign vectors on $n$ coordinates to signs,
let $P$ be a set of coordinates, and fix a sign assignment $z$ on $P$. Choose uniformly
a batch of $R$ sign vectors on all $n$ coordinates, then form the raw projected
restriction batch using $f$, $P$, $z$, and $M$. Its distribution equals the distribution
obtained by first choosing uniformly a batch of $R$ assignments on $P$, then
independently choosing uniformly a batch of $R$ assignments on the complementary
coordinates, combining each pair of assignments into a full sign vector, and applying
the same raw projection.
\end{theorem}

\begin{theorem}[Finiteness of a PMF outer measure]
\lean{FABL.pmfToOuterMeasure_ne_top}
\label{FABL.pmfToOuterMeasure_ne_top}
\leanok
\difficulty{3}
For any type $\alpha$, probability mass function $p$ on $\alpha$, and set
$E\subseteq\alpha$, the outer measure assigned to $E$ by $p$ is not infinite.
\end{theorem}

\begin{theorem}[Bad-event probability under junta projection]
\lean{FABL.rawProjectedJuntaRestrictionBatch_successBadProbability_le}
\label{FABL.rawProjectedJuntaRestrictionBatch_successBadProbability_le}
\leanok
\uses{FABL.BooleanFunction, FABL.JuntaFixedAssignment, FABL.JuntaFreeAssignment, FABL.Sign, FABL.combineJuntaAssignment, FABL.fixedMatchingIndices, FABL.juntaLabeledSamplesFromSplit, FABL.juntaMatchCount, FABL.juntaMatchingIndices, FABL.juntaRestrictionLabeledBatch, FABL.map_uniformPMF_rawProjectedJuntaRestrictionBatch, FABL.map_uniformPMF_selectedJuntaRestrictionBatch, FABL.matchedJuntaRestrictionExample, FABL.pmfToOuterMeasure_ne_top, FABL.projectJuntaRestrictionBatch, FABL.rawProjectedJuntaRestrictionBatch, FABL.selectedJuntaFreeBatch, FABL.takeMatchingJuntaExamples, FABL.takeMatchingJuntaExamples_eq_none_iff, FABL.takeMatchingJuntaExamples_from_split, FABL.uniformPMF}
\difficulty{6}
Let $n,M,R$ be nonnegative integers, let $f:\{-1,1\}^n\to\{\pm1\}$ be a sign-valued
Boolean function, let $P\subseteq\{1,\ldots,n\}$ be a set of coordinates, and fix a sign
assignment $z$ on $P$. Let $B$ be a property of batches of $M$ pairs, each consisting of
an assignment to the coordinates outside $P$ and a sign. If a uniformly random batch of
$M$ independent uniform assignments to those outside coordinates, labeled by the
restriction of $f$ with the fixed coordinates set to $z$, satisfies $B$ with probability
at most $\eta$, then for $R$ independent uniform inputs in $\{-1,1\}^n$, the probability
that the procedure selecting up to $M$ samples whose fixed coordinates match $z$ returns
a batch satisfying $B$ is at most $\eta$; failure to return a batch is not counted as a
bad outcome.
\end{theorem}

\begin{theorem}[Junta restriction sample program bad probability le]
\lean{FABL.juntaRestrictionSampleProgram_badProbability_le}
\statementsource{Analysis-of-Boolean-Functions-by-Ryan-ODonnell}{pdf-ffbe9c7b3306-p168-b001}
\label{FABL.juntaRestrictionSampleProgram_badProbability_le}
\leanok
\uses{FABL.BooleanFunction, FABL.JuntaFixedAssignment, FABL.JuntaFreeAssignment, FABL.JuntaRestrictionSampleBad, FABL.LearningCost, FABL.LearningProgram.eventProbability, FABL.LearningProgram.runWithCost, FABL.MatchedJuntaExample, FABL.PositiveLearningParameter, FABL.Sign, FABL.juntaRestrictionLabeledBatch, FABL.juntaRestrictionSampleCount, FABL.juntaRestrictionSampleProgram, FABL.juntaRestrictionSampleProgram_failureProbability_le, FABL.pmfToOuterMeasure_ne_top, FABL.projectJuntaRestrictionBatch, FABL.rawProjectedJuntaRestrictionBatch, FABL.rawProjectedJuntaRestrictionBatch_successBadProbability_le, FABL.runWithCost_juntaRestrictionSampleProgram, FABL.uniformPMF}
\difficulty{5}
% dataset-pipeline: stub -- prose pending build_informal_statements.py
Statement pending regeneration.
\end{theorem}

\begin{definition}[Relevant coordinate of a restricted function]
\lean{FABL.IsRelevantJuntaRestriction}
\statementsource{Analysis-of-Boolean-Functions-by-Ryan-ODonnell}{
  pdf-ffbe9c7b3306-p047-b005,
  pdf-ffbe9c7b3306-p075-b002,
  pdf-ffbe9c7b3306-p156-b004,
  pdf-ffbe9c7b3306-p168-b001,
  pdf-ffbe9c7b3306-p274-b007
}
\label{FABL.IsRelevantJuntaRestriction}
\leanok
\uses{FABL.BooleanFunction, FABL.JuntaFixedAssignment, FABL.JuntaFreeAssignment, FABL.JuntaFreeIndex, FABL.juntaRestriction}
Let $n$ be a natural number, let $f$ be a sign-valued Boolean function on the
$n$-dimensional sign cube, let $P$ be a set of coordinates, and fix signs on $P$. A
coordinate $i$ outside $P$ is relevant to the resulting restriction if there are two
sign assignments on the remaining coordinates that agree at every coordinate other than
$i$ but give different values of the restricted function.
\end{definition}

\begin{theorem}[Relevant restriction indices lie in the dependency set]
\lean{FABL.mem_of_isRelevantJuntaRestriction_of_dependsOn}
\label{FABL.mem_of_isRelevantJuntaRestriction_of_dependsOn}
\leanok
\uses{FABL.BooleanFunction, FABL.IsRelevantJuntaRestriction, FABL.JuntaFixedAssignment, FABL.JuntaFreeIndex, FABL.combineJuntaAssignment, FABL.juntaRestriction}
\difficulty{5}
Let $n$ be a natural number, let $f:\{\pm1\}^n\to\{\pm1\}$, and let $J,P$ be subsets of
the coordinate indices. Fix signs on $P$ and consider the restriction of $f$ obtained by
varying the remaining coordinates. If $f$ depends only on coordinates in $J$, and
changing the free coordinate $i\notin P$ can change this restricted function while all
other free coordinates are held fixed, then $i\in J$.
\end{theorem}

\begin{theorem}[Restriction is constant when all relevant coordinates are fixed]
\lean{FABL.juntaRestriction_eq_const_of_dependsOn_of_subset}
\label{FABL.juntaRestriction_eq_const_of_dependsOn_of_subset}
\leanok
\uses{FABL.BooleanFunction, FABL.JuntaFixedAssignment, FABL.combineJuntaAssignment, FABL.juntaRestriction}
\difficulty{4}
Let $n\in\mathbb{N}$, let $f:\{−1,1\}^n\to\{−1,1\}$ be a Boolean function, and let
$J,P\subseteq\{1,\ldots,n\}$. Fix an assignment $z$ on $P$. If $f$ depends only on
coordinates in $J$ and $J\subseteq P$, then the restriction of $f$ obtained by fixing
$P$ according to $z$ is constant on all assignments to the remaining coordinates, with
value given by evaluating $f$ at the cube point formed from $z$ on $P$ and the
all-positive assignment outside $P$.
\end{theorem}

\begin{definition}[Junta node decisions]
\lean{FABL.JuntaNodeDecision}
\statementsource{Analysis-of-Boolean-Functions-by-Ryan-ODonnell}{pdf-ffbe9c7b3306-p168-b001}
\label{FABL.JuntaNodeDecision}
\leanok
\uses{FABL.JuntaFreeIndex, FABL.Sign}
For a coordinate set $P$ among $n$ coordinates, a node decision is either to declare the
restricted function constant with a specified sign, or to split on a specified
coordinate outside $P$.
\end{definition}

\begin{definition}[Correct junta node decision]
\lean{FABL.JuntaNodeDecision.IsCorrect}
\label{FABL.JuntaNodeDecision.IsCorrect}
\leanok
\uses{FABL.BooleanFunction, FABL.IsRelevantJuntaRestriction, FABL.JuntaFixedAssignment, FABL.JuntaNodeDecision, FABL.juntaRestriction}
Let $n$ be a natural number, let $f$ be a Boolean function on the $n$-dimensional
Boolean cube, let $P$ be a finite set of coordinates in $	ext{Fin}(n)$, and let $z$
assign signs to the coordinates in $P$. A decision declaring a constant sign is correct
exactly when the restriction of $f$ obtained by fixing $P$ according to $z$ is
constantly that sign; a decision to split on a coordinate outside $P$ is correct exactly
when there are two assignments to the remaining coordinates that agree everywhere except
at that coordinate and give different values of the restriction.
\end{definition}

\begin{definition}[Bad junta node decision]
\lean{FABL.JuntaNodeDecision.IsBad}
\label{FABL.JuntaNodeDecision.IsBad}
\leanok
\uses{FABL.BooleanFunction, FABL.JuntaFixedAssignment, FABL.JuntaNodeDecision, FABL.JuntaNodeDecision.IsCorrect}
Let $n$ be a natural number, let $f$ be a Boolean function on the $n$-dimensional
Boolean cube, let $P$ be a finite set of coordinates in $	ext{Fin}(n)$, and let $z$
assign signs to the coordinates in $P$. An optional node decision is bad if it is
absent, or, when present, if it is not correct for the restriction of $f$ determined by
$P$ and $z$.
\end{definition}

\begin{theorem}[Correctness of a nonbad decision]
\lean{FABL.JuntaNodeDecision.isCorrect_of_not_isBad}
\label{FABL.JuntaNodeDecision.isCorrect_of_not_isBad}
\leanok
\uses{FABL.BooleanFunction, FABL.JuntaFixedAssignment, FABL.JuntaNodeDecision, FABL.JuntaNodeDecision.IsBad, FABL.JuntaNodeDecision.IsCorrect}
\difficulty{1}
Let $n\ge 0$, let $P$ be a set of coordinates among the $n$ inputs, and fix signs on the
coordinates in $P$. Let the target be a function from the sign cube ${-1,1}^n$ to
$	ext{signs }{pm1}$. A decision is either that the target restricted to the remaining
coordinates is constant at a specified sign, or that it depends on a specified remaining
coordinate. If this decision is not bad—that is, it is not incorrect—then it is correct:
in the constant case the restriction equals the specified constant function, and in the
coordinate case the restriction genuinely depends on that coordinate.
\end{theorem}

\begin{definition}[Relevant coordinate finder]
\lean{FABL.JuntaRelevantCoordinateFinder}
\statementsource{Analysis-of-Boolean-Functions-by-Ryan-ODonnell}{
  pdf-ffbe9c7b3306-p089-b008,
  pdf-ffbe9c7b3306-p156-b004,
  pdf-ffbe9c7b3306-p156-b005,
  pdf-ffbe9c7b3306-p168-b001
}
\label{FABL.JuntaRelevantCoordinateFinder}
\leanok
\uses{FABL.BooleanFunction, FABL.JuntaFixedAssignment, FABL.JuntaNodeDecision, FABL.JuntaNodeDecision.IsBad, FABL.LearningCost, FABL.LearningProgram, FABL.LearningProgram.eventProbability, FABL.LearningProgram.runWithCost, FABL.PositiveLearningParameter}
For natural numbers $n$ and $k$, a relevant coordinate finder assigns to every
coordinate set $P$, fixed sign assignment on $P$, and rational failure parameter
strictly between $0$ and $1/2$ a randomized-example learning program returning an
optional node decision. It also specifies bounds, depending on the failure parameter, on
the number of random examples and local-work units; every supported outcome must use at
most those bounds and make zero membership queries. For every sign-valued Boolean
function, every coordinate sets $J,P$ with $|J|\le k$, every fixed assignment on $P$,
and every inclusion $P\subseteq J$, if the function depends only on coordinates in $J$,
then the probability that the program returns no decision or an incorrect decision is at
most the failure parameter.
\end{definition}

\begin{definition}[Binary tree call count]
\lean{FABL.juntaTreeCallCount}
\label{FABL.juntaTreeCallCount}
\leanok
The number of calls in a full binary tree of depth parameter $k$ is defined recursively
by $C(0)=1$ and $C(k+1)=1+2C(k)$.
\end{definition}

\begin{theorem}[Initial junta tree call count]
\lean{FABL.juntaTreeCallCount_zero}
\label{FABL.juntaTreeCallCount_zero}
\leanok
\uses{FABL.juntaTreeCallCount}
\difficulty{1}
The junta tree call count at depth zero is one.
\end{theorem}

\begin{theorem}[Successor recurrence for the call count]
\lean{FABL.juntaTreeCallCount_succ}
\label{FABL.juntaTreeCallCount_succ}
\leanok
\uses{FABL.juntaTreeCallCount}
\difficulty{1}
For every nonnegative integer $k$, the junta tree call count at depth $k+1$ equals one
plus twice the call count at depth $k$.
\end{theorem}

\begin{theorem}[Positivity of the junta tree call count]
\lean{FABL.juntaTreeCallCount_pos}
\label{FABL.juntaTreeCallCount_pos}
\leanok
\uses{FABL.juntaTreeCallCount}
\difficulty{4}
For every nonnegative integer $k$, the junta tree call count at depth $k$ is positive.
\end{theorem}

\begin{definition}[Per-call failure parameter]
\lean{FABL.juntaTreePerCallFailure}
\label{FABL.juntaTreePerCallFailure}
\leanok
\uses{FABL.PositiveLearningParameter, FABL.juntaTreeCallCount, FABL.juntaTreeCallCount_pos}
For a natural number $k$ and a rational failure parameter $ailure\in(0,1/2]$, define
the per-call failure parameter to be $ailure/C(k)$, where $C$ is the binary tree call
count.
\end{definition}

\begin{theorem}[Per-call failure parameter formula]
\lean{FABL.juntaTreePerCallFailure_value}
\label{FABL.juntaTreePerCallFailure_value}
\leanok
\uses{FABL.PositiveLearningParameter, FABL.juntaTreeCallCount, FABL.juntaTreePerCallFailure}
\difficulty{1}
Let $k$ be a nonnegative integer and let $	au$ be a rational learning parameter with
$0<	au\le\tfrac12$. The per-call failure parameter at depth $k$ is $	au$ divided by the
junta tree call count at depth $k$.
\end{theorem}

\begin{theorem}[Junta tree failure allocation]
\lean{FABL.juntaTreeCallCount_mul_perCallFailure}
\label{FABL.juntaTreeCallCount_mul_perCallFailure}
\leanok
\uses{FABL.PositiveLearningParameter, FABL.juntaTreeCallCount, FABL.juntaTreeCallCount_pos, FABL.juntaTreePerCallFailure, FABL.juntaTreePerCallFailure_value}
\difficulty{3}
Let $k$ be a nonnegative integer and let $	au$ be a rational learning parameter with
$0<	au\le \tfrac12$. If the per-call failure parameter is obtained by dividing $	au$ by
the junta tree call count at depth $k$, then the call count multiplied by that per-call
parameter is exactly $	au$.
\end{theorem}

\begin{definition}[Insert a fixed assignment]
\lean{FABL.insertJuntaFixedAssignment}
\label{FABL.insertJuntaFixedAssignment}
\leanok
\uses{FABL.JuntaFixedAssignment, FABL.JuntaFreeIndex, FABL.Sign}
Let $P$ be a subset of the coordinates $in n$, let $z$ assign a sign in $\{,-1\}$ to
each coordinate in $P$, and let $i$ be a coordinate outside $P$. For a sign $s$, the
assignment on $P\cup\{i\}$ obtained by inserting $i$ assigns $s$ to $i$ and agrees with
$z$ on every coordinate in $P$.
\end{definition}

\begin{theorem}[Value at the newly inserted coordinate]
\lean{FABL.insertJuntaFixedAssignment_apply_new}
\label{FABL.insertJuntaFixedAssignment_apply_new}
\leanok
\uses{FABL.JuntaFixedAssignment, FABL.JuntaFreeIndex, FABL.Sign, FABL.insertJuntaFixedAssignment}
\difficulty{1}
Let $P$ be a set of coordinates in $\mathrm{Fin}(n)$, let $z$ assign a sign to each
coordinate in $P$, and let $i\notin P$ be a coordinate and $v$ a sign. Extend the
assignment by inserting $i$ with value $v$ while retaining the old values on $P$. Its
value at $i$ is $v$.
\end{theorem}

\begin{theorem}[Values retained on the original coordinates]
\lean{FABL.insertJuntaFixedAssignment_apply_old}
\label{FABL.insertJuntaFixedAssignment_apply_old}
\leanok
\uses{FABL.JuntaFixedAssignment, FABL.JuntaFreeIndex, FABL.Sign, FABL.insertJuntaFixedAssignment}
\difficulty{4}
Let $P$ be a set of coordinates in $\mathrm{Fin}(n)$, let $z$ assign a sign to each
coordinate in $P$, and let $i\notin P$ be a coordinate and $v$ a sign. Extend the
assignment by inserting $i$ with value $v$ while retaining the old values on $P$. At
every coordinate in $P$, the extended assignment agrees with $z$.
\end{theorem}

\begin{theorem}[Matching an inserted fixed assignment]
\lean{FABL.matches_insertJuntaFixedAssignment_iff}
\label{FABL.matches_insertJuntaFixedAssignment_iff}
\leanok
\uses{FABL.JuntaFixedAssignment, FABL.JuntaFreeIndex, FABL.MatchesJuntaAssignment, FABL.Sign, FABL.insertJuntaFixedAssignment}
\difficulty{5}
Let $n$ be a natural number, let $P\\subseteq\\{1,\\ldots,n\\}$, and let $z$ assign
signs to the coordinates in $P$. Choose a coordinate $i\\notin P$, a sign $v$, and an
input $x\\in\\{-1,1\\}^n$. Then $x$ agrees with the fixed assignment obtained by
extending $z$ to set coordinate $i$ to $v$ if and only if $x$ agrees with $z$ on every
coordinate in $P$ and $x_i=v$.
\end{theorem}

\begin{definition}[Cast a decision tree's available coordinates]
\lean{FABL.F₂DecisionTree.castAvailable}
\leanok
\uses{FABL.F₂DecisionTree}
Given finite coordinate sets $A$ and $B$ with $A=B$, and a binary decision tree whose
available coordinates are $A$, transport the tree along this equality to view it as a
tree whose available coordinates are $B$.
\end{definition}

\begin{theorem}[Evaluation after changing the available set]
\lean{FABL.F₂DecisionTree.eval_castAvailable}
\leanok
\uses{FABL.F₂Cube, FABL.F₂DecisionTree, FABL.F₂DecisionTree.castAvailable, FABL.F₂DecisionTree.eval}
\difficulty{1}
Let $A$ and $B$ be finite sets of coordinates with $A=B$, and let $T$ be a binary
decision tree on $A$. For every input in the $n$-dimensional binary cube, regarding $T$
as a tree on $B$ by this equality gives the same output as evaluating $T$ on that input.
\end{theorem}

\begin{theorem}[Depth invariance under reindexing availability]
\lean{FABL.F₂DecisionTree.depth_castAvailable}
\leanok
\uses{FABL.F₂DecisionTree, FABL.F₂DecisionTree.castAvailable, FABL.F₂DecisionTree.depth}
\difficulty{1}
Let $n$ be a nonnegative integer, let $A$ and $B$ be sets of coordinates with $A=B$, and
let $T$ be a decision tree indexed by the available coordinate set $A$. Reindexing $T$
as a tree with available coordinate set $B$ leaves its depth unchanged.
\end{theorem}

\begin{definition}[Assembling a junta query]
\lean{FABL.assembleJuntaQuery}
\statementsource{Analysis-of-Boolean-Functions-by-Ryan-ODonnell}{pdf-ffbe9c7b3306-p168-b001}
\label{FABL.assembleJuntaQuery}
\leanok
\uses{FABL.F₂DecisionTree, FABL.F₂DecisionTree.castAvailable, FABL.JuntaFreeIndex, FABL.Sign}
Let $P$ be a set of coordinates among $n$, and let $i$ be a coordinate outside $P$.
Given optional decision trees for the two branches, each with available coordinates
outside $P\cup\{i\}$, return the tree that queries coordinate $i$ and uses those trees
as its two children if both are present; otherwise return no tree.
\end{definition}

\begin{definition}[Leaf output from a node decision]
\lean{FABL.juntaNodeLeafOutput}
\statementsource{Analysis-of-Boolean-Functions-by-Ryan-ODonnell}{pdf-ffbe9c7b3306-p168-b001}
\label{FABL.juntaNodeLeafOutput}
\leanok
\uses{FABL.F₂DecisionTree, FABL.JuntaNodeDecision, FABL.Sign}
For a coordinate set $P$ among $n$ coordinates, an optional node decision is sent to an
optional decision tree with available coordinates outside $P$: a present constant
decision becomes a leaf storing that sign, and every other decision becomes no tree.
\end{definition}

\begin{definition}[Recursive junta learner construction]
\lean{FABL.recursiveJuntaLearnerAux}
\statementsource{Analysis-of-Boolean-Functions-by-Ryan-ODonnell}{
  pdf-ffbe9c7b3306-p156-b004,
  pdf-ffbe9c7b3306-p156-b005,
  pdf-ffbe9c7b3306-p156-b006,
  pdf-ffbe9c7b3306-p156-b007,
  pdf-ffbe9c7b3306-p156-b008,
  pdf-ffbe9c7b3306-p168-b001
}
\label{FABL.recursiveJuntaLearnerAux}
\leanok
\uses{FABL.F₂DecisionTree, FABL.JuntaFixedAssignment, FABL.JuntaRelevantCoordinateFinder, FABL.LearningProgram, FABL.LearningProgram.bind, FABL.LearningProgram.map, FABL.PositiveLearningParameter, FABL.Sign, FABL.assembleJuntaQuery, FABL.binarySignEquiv, FABL.insertJuntaFixedAssignment, FABL.juntaNodeLeafOutput}
Let $n$ and $k$ be natural numbers, let a relevant-coordinate finder be given for
dimension $n$ and parameter $k$, and let the rational failure parameter lie in
$(0,1/2]$. For every natural number of fuel, finite coordinate set
$P\subseteq\mathrm{Fin}(n)$, and fixed sign assignment on $P$, the construction produces
a randomized-example learning program returning either no decision tree or a decision
tree whose available coordinates are outside $P$. With zero fuel, it maps the finder’s
output to a leaf only when the output declares a constant; with positive fuel, it
returns no tree for an absent decision, a leaf for a constant decision, and, for a
relevant coordinate, recursively constructs the two branches with that coordinate fixed
to each of the two signs and assembles them into a query tree.
\end{definition}

\begin{definition}[Empty fixed assignment]
\lean{FABL.emptyJuntaFixedAssignment}
\label{FABL.emptyJuntaFixedAssignment}
\leanok
\uses{FABL.JuntaFixedAssignment}
For any natural number $n$, the fixed assignment on the empty subset of $in n$ is the
unique function from that empty set to the signs.
\end{definition}

\begin{definition}[Recursive junta learner]
\lean{FABL.recursiveJuntaLearner}
\statementsource{Analysis-of-Boolean-Functions-by-Ryan-ODonnell}{
  pdf-ffbe9c7b3306-p156-b002,
  pdf-ffbe9c7b3306-p156-b004,
  pdf-ffbe9c7b3306-p156-b005,
  pdf-ffbe9c7b3306-p156-b006,
  pdf-ffbe9c7b3306-p156-b007,
  pdf-ffbe9c7b3306-p156-b008,
  pdf-ffbe9c7b3306-p168-b001
}
\label{FABL.recursiveJuntaLearner}
\leanok
\uses{FABL.DecisionTree, FABL.F₂DecisionTree.castAvailable, FABL.JuntaRelevantCoordinateFinder, FABL.LearningProgram, FABL.LearningProgram.map, FABL.PositiveLearningParameter, FABL.Sign, FABL.emptyJuntaFixedAssignment, FABL.juntaTreePerCallFailure, FABL.recursiveJuntaLearnerAux}
Let $n,k$ be natural numbers, let a relevant coordinate finder be given for these
parameters, and choose a rational failure parameter in $(0,1/2]$. The learner is the
randomized-example program obtained by running the recursive construction with fuel $k$,
starting with no fixed coordinates and allocating the per-call failure parameter
prescribed by the learner’s failure-allocation rule; its output, when present, is a
decision tree on all $n$ coordinates.
\end{definition}

\begin{definition}[Decision tree computes a junta restriction]
\lean{FABL.F₂DecisionTree.ComputesJuntaRestriction}
\leanok
\uses{FABL.BooleanFunction, FABL.F₂Cube, FABL.F₂DecisionTree, FABL.F₂DecisionTree.eval, FABL.JuntaFixedAssignment, FABL.MatchesJuntaAssignment, FABL.Sign, FABL.binaryCubeSignEquiv}
\statementsource{Analysis-of-Boolean-Functions-by-Ryan-ODonnell}{
  pdf-ffbe9c7b3306-p073-b004,
  pdf-ffbe9c7b3306-p075-b002,
  pdf-ffbe9c7b3306-p168-b001
}
Let $n$ be a natural number, let $P$ be a set of coordinates in $\\{0,\ldots,n-1\}$, let
$z$ assign a sign in $\{ -1,1 \}$ to each coordinate in $P$, and let $T$ be a decision
tree whose available query coordinates are precisely those outside $P$. Then $T$
computes the target function on every Boolean-cube input whose signs on $P$ agree with
$z$.
\end{definition}

\begin{theorem}[Leaf computes a constant junta restriction]
\lean{FABL.F₂DecisionTree.computesJuntaRestriction_leaf}
\leanok
\uses{FABL.BooleanFunction, FABL.F₂DecisionTree, FABL.F₂DecisionTree.ComputesJuntaRestriction, FABL.JuntaFixedAssignment, FABL.Sign, FABL.binaryCubeSignEquiv, FABL.combineJuntaAssignment_freePart_of_matches, FABL.juntaFreePart, FABL.juntaRestriction}
\difficulty{3}
Let $n$ be a natural number, let $f$ map the sign cube $\{\pm1\}^n$ to the two signs,
let $P$ be a set of coordinates, and fix a sign assignment $z$ on $P$. If the
restriction of $f$ obtained by fixing the coordinates in $P$ according to $z$ is
constantly equal to a sign $v$, then the decision tree consisting of a single leaf
labelled $v$, with available coordinates those outside $P$, computes $f$ on every binary
input whose sign encoding agrees with $z$ on $P$.
\end{theorem}

\begin{definition}[Bad restricted tree output]
\lean{FABL.JuntaTreeOutputBad}
\label{FABL.JuntaTreeOutputBad}
\leanok
\uses{FABL.BooleanFunction, FABL.F₂DecisionTree, FABL.F₂DecisionTree.ComputesJuntaRestriction, FABL.JuntaFixedAssignment, FABL.Sign}
Let $f$ be a Boolean function on the $n$-dimensional sign cube, let $Pinin n$, and let
$z$ assign signs to the coordinates in $P$. An absent tree output is declared not bad; a
present tree output is bad exactly when the tree fails to compute $f$ on every input
whose coordinates in $P$ agree with $z$, using only coordinates outside $P$ for its
queries.
\end{definition}

\begin{theorem}[Decision tree computes the junta restriction]
\lean{FABL.computesJuntaRestriction_of_not_outputBad}
\label{FABL.computesJuntaRestriction_of_not_outputBad}
\leanok
\uses{FABL.BooleanFunction, FABL.F₂DecisionTree, FABL.F₂DecisionTree.ComputesJuntaRestriction, FABL.JuntaFixedAssignment, FABL.JuntaTreeOutputBad, FABL.Sign}
\difficulty{1}
Let $n$ be a natural number, let $f$ be a sign-valued Boolean function on $n$
coordinates, let $P$ be a set of coordinates, and fix a sign assignment $z$ on $P$. If a
decision tree querying only coordinates outside $P$ is not classified as a bad output
for $(f,P,z)$, then on every binary input whose sign vector agrees with $z$ on $P$, the
tree returns the value of $f$ at that sign vector.
\end{theorem}

\begin{theorem}[Combining successful junta restrictions]
\lean{FABL.not_outputBad_assembleJuntaQuery}
\label{FABL.not_outputBad_assembleJuntaQuery}
\leanok
\uses{FABL.BooleanFunction, FABL.F₂DecisionTree, FABL.F₂DecisionTree.ComputesJuntaRestriction, FABL.F₂DecisionTree.castAvailable, FABL.F₂DecisionTree.eval, FABL.F₂DecisionTree.eval_castAvailable, FABL.JuntaFixedAssignment, FABL.JuntaFreeIndex, FABL.JuntaTreeOutputBad, FABL.Sign, FABL.assembleJuntaQuery, FABL.binaryCubeSignEquiv_apply, FABL.binarySignEquiv, FABL.computesJuntaRestriction_of_not_outputBad, FABL.insertJuntaFixedAssignment, FABL.matches_insertJuntaFixedAssignment_iff}
\difficulty{5}
Let $f$ be a function from the $n$-dimensional sign cube to $\{\pm1\}$, let $P$ be a set
of coordinates, let $z$ assign a sign to each coordinate in $P$, and let $i$ be a
coordinate outside $P$. Suppose each of two optional decision trees, with coordinates
available outside $P\cup\{i\}$, is either absent or correctly computes the restriction
of $f$ for the assignment extending $z$ by setting coordinate $i$ to the sign encoded by
$0$ or $1$, respectively. Then the optional tree obtained by querying coordinate $i$ and
using these two trees as its branches is either absent or correctly computes the
restriction of $f$ for the assignment $z$.
\end{theorem}

\begin{theorem}[Subset from zero set difference]
\lean{FABL.subset_of_card_sdiff_le_zero}
\label{FABL.subset_of_card_sdiff_le_zero}
\leanok
\difficulty{4}
Let $n$ be a natural number and let $J,P\subseteq\{0,\ldots,n-1\}$. If the set
difference $J\setminus P$ has cardinality at most zero, then $J\subseteq P$.
\end{theorem}

\begin{theorem}[Cardinality after removing an inserted index]
\lean{FABL.card_sdiff_insert_le_of_mem_of_notMem}
\label{FABL.card_sdiff_insert_le_of_mem_of_notMem}
\leanok
\difficulty{3}
Let $J$ and $P$ be subsets of the $n$ coordinate indices, let $i$ be a coordinate, and
let $f$ be a natural number. If $i\in J$, $i\notin P$, and $|J\setminus P|\le f+1$, then
$|J\setminus(P\cup\{i\})|\le f$.
\end{theorem}

\begin{theorem}[Correct junta decision with no free relevant coordinates]
\lean{FABL.exists_constant_of_correctNode_of_card_sdiff_le_zero}
\label{FABL.exists_constant_of_correctNode_of_card_sdiff_le_zero}
\leanok
\uses{FABL.BooleanFunction, FABL.JuntaFixedAssignment, FABL.JuntaNodeDecision, FABL.JuntaNodeDecision.IsCorrect, FABL.Sign, FABL.mem_of_isRelevantJuntaRestriction_of_dependsOn, FABL.subset_of_card_sdiff_le_zero}
\difficulty{4}
Let $f$ be a sign-valued function on the Boolean cube, let $J$ and $P$ be sets of
coordinates, and fix values on the coordinates in $P$. Suppose that $f$ depends only on
coordinates in $J$, that $|J\setminus P|\le 0$, and that a decision for the restriction
of $f$ to the fixed assignment is correct. Then the decision is a constant decision:
there is a sign value that it returns for every assignment to the remaining coordinates.
\end{theorem}

\begin{theorem}[Output maps preserve event probabilities]
\lean{FABL.eventProbability_map_output_eq}
\label{FABL.eventProbability_map_output_eq}
\leanok
\uses{FABL.BooleanFunction, FABL.LearningAccess, FABL.LearningCost, FABL.LearningProgram, FABL.LearningProgram.eventProbability, FABL.LearningProgram.map, FABL.LearningProgram.runWithCost_map}
\difficulty{3}
Let $n$ be a natural number, let a learning program on the $n$-dimensional sign cube
have output space $A$, let $f:A\to B$ be any map, and let $T$ be a Boolean function on
that cube. For every predicate $Q$ on $B$, the probability that the program obtained by
applying $f$ to the output has output satisfying $Q$ equals the probability that the
original program has output $a$ satisfying $Q(f(a))$.
\end{theorem}

\begin{theorem}[Output maps bound bad-event probabilities]
\lean{FABL.eventProbability_map_output_le}
\label{FABL.eventProbability_map_output_le}
\leanok
\uses{FABL.BooleanFunction, FABL.LearningAccess, FABL.LearningCost, FABL.LearningProgram, FABL.LearningProgram.eventProbability, FABL.LearningProgram.eventProbability_mono, FABL.LearningProgram.map, FABL.eventProbability_map_output_eq}
\difficulty{3}
Let $n$ be a natural number, let a learning program on the $n$-dimensional sign cube
have output space $A$, let $f:A\to B$ be a map, and let $T$ be a Boolean function on
that cube. Suppose $P$ is a predicate on $A$, $Q$ is a predicate on $B$, and $Q(f(a))$
implies $P(a)$ for every $a\in A$. Then the probability that the mapped program's output
satisfies $Q$ is at most the probability that the original program's output satisfies
$P$.
\end{theorem}

\begin{theorem}[A good pure output has zero bad-event probability]
\lean{FABL.eventProbability_pure_eq_zero_of_not}
\label{FABL.eventProbability_pure_eq_zero_of_not}
\leanok
\uses{FABL.BooleanFunction, FABL.LearningAccess, FABL.LearningProgram, FABL.LearningProgram.eventProbability, FABL.LearningProgram.runWithCost}
\difficulty{1}
Let $n$ be a natural number, let $T$ be a Boolean function on the $n$-dimensional sign
cube, and let a learning program return a fixed value $a$ without random choices. If a
predicate $P$ on its output space does not hold for $a$, then the probability that the
output satisfies $P$ is zero.
\end{theorem}

\begin{theorem}[Bad leaf output forces a bad decision]
\lean{FABL.nodeLeafOutput_bad_implies_nodeBad}
\label{FABL.nodeLeafOutput_bad_implies_nodeBad}
\leanok
\uses{FABL.BooleanFunction, FABL.F₂DecisionTree.computesJuntaRestriction_leaf, FABL.JuntaFixedAssignment, FABL.JuntaNodeDecision, FABL.JuntaNodeDecision.IsBad, FABL.JuntaNodeDecision.isCorrect_of_not_isBad, FABL.JuntaTreeOutputBad, FABL.exists_constant_of_correctNode_of_card_sdiff_le_zero, FABL.juntaNodeLeafOutput}
\difficulty{4}
Let $n$ be a natural number, let $\mathrm{target}$ be a sign-valued function on the
$n$-dimensional sign cube, and let $J,P\subseteq\{1,\ldots,n\}$. Fix signs on the
coordinates in $P$, and suppose the target is determined by the coordinates in $J$ and
$|J\setminus P|\le0$. For any node decision, if the output obtained by replacing a
constant decision with its corresponding leaf (and any other decision with no output)
fails to compute the target on inputs agreeing with the fixed signs, then the node
decision itself is incorrect (with a missing decision counted as bad).
\end{theorem}

\begin{theorem}[Failure probability of the recursive learner]
\lean{FABL.recursiveJuntaLearnerAux_failureProbability_le}
\proofsource{Analysis-of-Boolean-Functions-by-Ryan-ODonnell}{
  pdf-ffbe9c7b3306-p156-b006,
  pdf-ffbe9c7b3306-p168-b001
}
\label{FABL.recursiveJuntaLearnerAux_failureProbability_le}
\leanok
\uses{FABL.BooleanFunction, FABL.F₂DecisionTree.computesJuntaRestriction_leaf, FABL.JuntaFixedAssignment, FABL.JuntaNodeDecision.IsBad, FABL.JuntaNodeDecision.isCorrect_of_not_isBad, FABL.JuntaRelevantCoordinateFinder, FABL.JuntaTreeOutputBad, FABL.LearningProgram.bind, FABL.LearningProgram.eventProbability, FABL.LearningProgram.eventProbability_bind_le_add, FABL.LearningProgram.map, FABL.PositiveLearningParameter, FABL.assembleJuntaQuery, FABL.binarySignEquiv, FABL.card_sdiff_insert_le_of_mem_of_notMem, FABL.eventProbability_map_output_le, FABL.eventProbability_pure_eq_zero_of_not, FABL.insertJuntaFixedAssignment, FABL.juntaNodeLeafOutput, FABL.juntaTreeCallCount, FABL.juntaTreeCallCount_succ, FABL.mem_of_isRelevantJuntaRestriction_of_dependsOn, FABL.nodeLeafOutput_bad_implies_nodeBad, FABL.not_outputBad_assembleJuntaQuery, FABL.recursiveJuntaLearnerAux}
\difficulty{6}
Let $n,k$ be nonnegative integers, let a node finder for $k$-juntas on the
$n$-dimensional Boolean cube be given, and let $f$ be a Boolean function on that cube
taking values in $\{\pm1\}$. Suppose $J$ is a set of coordinates with $|J|\le k$ on
which $f$ depends, and fix a learning parameter $\varepsilon\in(0,1/2]$. For every fuel
$r$, fixed-coordinate set $P\subseteq J$, and sign assignment on $P$ with $|J\setminus
P|\le r$, the probability that the recursive learner fails to return a decision tree
computing the restriction of $f$ consistent with that assignment is at most
$T(r)\varepsilon$, where $T(0)=1$ and $T(r+1)=1+2T(r)$.
\end{theorem}

\begin{theorem}[Depth bound for the recursive learner]
\lean{FABL.recursiveJuntaLearnerAux_depth_le}
\label{FABL.recursiveJuntaLearnerAux_depth_le}
\leanok
\uses{FABL.BooleanFunction, FABL.F₂DecisionTree, FABL.F₂DecisionTree.castAvailable, FABL.F₂DecisionTree.depth, FABL.F₂DecisionTree.depth_castAvailable, FABL.JuntaFixedAssignment, FABL.JuntaRelevantCoordinateFinder, FABL.LearningCost, FABL.LearningProgram.addOutcomeCost, FABL.LearningProgram.runWithCost, FABL.LearningProgram.runWithCost_bind, FABL.LearningProgram.runWithCost_map, FABL.PositiveLearningParameter, FABL.Sign, FABL.assembleJuntaQuery, FABL.binarySignEquiv, FABL.insertJuntaFixedAssignment, FABL.juntaNodeLeafOutput, FABL.recursiveJuntaLearnerAux}
\difficulty{6}
Let $n,k$ be nonnegative integers, let a node finder for $k$-juntas on the
$n$-dimensional Boolean cube be given, let $f$ be a Boolean function on that cube taking
values in $\{\pm1\}$, and fix a learning parameter $\varepsilon\in(0,1/2]$. For any fuel
$r$, set $P$ of fixed coordinates, and assignment of signs on $P$, every supported
outcome of the recursive learner that returns a decision tree has a tree of depth at
most $r$.
\end{theorem}

\begin{theorem}[Projected cost of the recursive learner]
\lean{FABL.recursiveJuntaLearnerAux_costProjection_le}
\label{FABL.recursiveJuntaLearnerAux_costProjection_le}
\leanok
\uses{FABL.BooleanFunction, FABL.F₂DecisionTree, FABL.JuntaFixedAssignment, FABL.JuntaNodeDecision, FABL.JuntaRelevantCoordinateFinder, FABL.LearningCost, FABL.LearningProgram.addOutcomeCost, FABL.LearningProgram.runWithCost, FABL.LearningProgram.runWithCost_bind, FABL.LearningProgram.runWithCost_map, FABL.PositiveLearningParameter, FABL.Sign, FABL.binarySignEquiv, FABL.insertJuntaFixedAssignment, FABL.juntaTreeCallCount, FABL.juntaTreeCallCount_succ, FABL.recursiveJuntaLearnerAux}
\difficulty{5}
Let $n,k$ be nonnegative integers, let a node finder for $k$-juntas on the
$n$-dimensional Boolean cube be given, and let $f$ be a Boolean function on that cube
taking values in $\{\pm1\}$. Fix a learning parameter $\varepsilon\in(0,1/2]$, an
additive map $q$ from learning costs to nonnegative integers with $q(0)=0$, and a
nonnegative integer $B$. Suppose every outcome in the support of each node-finder call,
for every set $P$ of fixed coordinates and every assignment of signs on $P$, has
projected cost at most $B$. Then for every recursion fuel $r$, every such $P$ and
assignment, and every outcome in the support of the recursive learner, its projected
cost is at most $T(r)B$, where $T(0)=1$ and $T(r+1)=1+2T(r)$.
\end{theorem}

\begin{theorem}[Recursive junta learner aux random examples le]
\lean{FABL.recursiveJuntaLearnerAux_randomExamples_le}
\statementsource{Analysis-of-Boolean-Functions-by-Ryan-ODonnell}{
  pdf-ffbe9c7b3306-p156-b006,
  pdf-ffbe9c7b3306-p156-b008,
  pdf-ffbe9c7b3306-p168-b001
}
\label{FABL.recursiveJuntaLearnerAux_randomExamples_le}
\leanok
\uses{FABL.BooleanFunction, FABL.F₂DecisionTree, FABL.JuntaFixedAssignment, FABL.JuntaRelevantCoordinateFinder, FABL.LearningCost, FABL.LearningProgram.runWithCost, FABL.PositiveLearningParameter, FABL.Sign, FABL.juntaTreeCallCount, FABL.recursiveJuntaLearnerAux, FABL.recursiveJuntaLearnerAux_costProjection_le}
\difficulty{3}
% dataset-pipeline: stub -- prose pending build_informal_statements.py
Statement pending regeneration.
\end{theorem}

\begin{theorem}[Zero query cost]
\lean{FABL.recursiveJuntaLearnerAux_queries_eq_zero}
\label{FABL.recursiveJuntaLearnerAux_queries_eq_zero}
\leanok
\uses{FABL.BooleanFunction, FABL.F₂DecisionTree, FABL.JuntaFixedAssignment, FABL.JuntaRelevantCoordinateFinder, FABL.LearningCost, FABL.LearningProgram.runWithCost, FABL.PositiveLearningParameter, FABL.Sign, FABL.recursiveJuntaLearnerAux, FABL.recursiveJuntaLearnerAux_costProjection_le}
\difficulty{3}
Let $n,k$ be nonnegative integers, let a node finder for $k$-juntas on the
$n$-dimensional Boolean cube be given, let $f$ be a Boolean function on that cube taking
values in $\{\pm1\}$, and fix a learning parameter $\varepsilon\in(0,1/2]$. For every
fuel, set $P$ of fixed coordinates, assignment of signs on $P$, and outcome in the
support of the recursive learner's output-and-cost distribution, the number of
membership queries charged to that outcome is zero.
\end{theorem}

\begin{theorem}[Work bound for the recursive learner]
\lean{FABL.recursiveJuntaLearnerAux_work_le}
\label{FABL.recursiveJuntaLearnerAux_work_le}
\leanok
\uses{FABL.BooleanFunction, FABL.F₂DecisionTree, FABL.JuntaFixedAssignment, FABL.JuntaRelevantCoordinateFinder, FABL.LearningCost, FABL.LearningProgram.runWithCost, FABL.PositiveLearningParameter, FABL.Sign, FABL.juntaTreeCallCount, FABL.recursiveJuntaLearnerAux, FABL.recursiveJuntaLearnerAux_costProjection_le}
\difficulty{3}
Let $n,k$ be nonnegative integers, let a node finder for $k$-juntas on the
$n$-dimensional Boolean cube be given, let $f$ be a Boolean function on that cube taking
values in $\{\pm1\}$, and fix a learning parameter $\varepsilon\in(0,1/2]$. For every
fuel $r$, set $P$ of fixed coordinates, assignment of signs on $P$, and outcome in the
support of the recursive learner's output-and-cost distribution, the charged local work
is at most $T(r)W(\varepsilon)$, where $W(\varepsilon)$ is the node finder's work bound
and $T(0)=1$, $T(r+1)=1+2T(r)$.
\end{theorem}

\begin{definition}[Decision tree computes a target]
\lean{FABL.DecisionTreeComputesTarget}
\statementsource{Analysis-of-Boolean-Functions-by-Ryan-ODonnell}{
  pdf-ffbe9c7b3306-p073-b004,
  pdf-ffbe9c7b3306-p074-b004,
  pdf-ffbe9c7b3306-p168-b001
}
\label{FABL.DecisionTreeComputesTarget}
\leanok
\uses{FABL.BooleanFunction, FABL.DecisionTree, FABL.F₂Cube, FABL.F₂DecisionTree.eval, FABL.Sign, FABL.binaryCubeSignEquiv}
Let $n$ be a natural number, let $T$ be a decision tree on all $n$ coordinates with
sign-valued output, and let $f$ be a sign-valued function on the $n$-dimensional Boolean
cube. Then for every input $x$ in the $_2$ cube, $T(x)=f(s(x))$, where $s(x)$ is the
corresponding sign vector.
\end{definition}

\begin{definition}[Bad learner output]
\lean{FABL.JuntaLearnerOutputBad}
\label{FABL.JuntaLearnerOutputBad}
\leanok
\uses{FABL.BooleanFunction, FABL.DecisionTree, FABL.DecisionTreeComputesTarget, FABL.Sign}
Let $f$ be a Boolean function on the $n$-dimensional sign cube. An absent decision-tree
output is declared not bad; a present output is bad exactly when the tree does not
compute $f$ on all inputs in the cube.
\end{definition}

\begin{theorem}[Bad root output gives bad auxiliary output]
\lean{FABL.rootOutput_bad_implies_auxBad}
\label{FABL.rootOutput_bad_implies_auxBad}
\leanok
\uses{FABL.BooleanFunction, FABL.F₂DecisionTree, FABL.F₂DecisionTree.castAvailable, FABL.F₂DecisionTree.eval_castAvailable, FABL.JuntaLearnerOutputBad, FABL.JuntaTreeOutputBad, FABL.Sign, FABL.emptyJuntaFixedAssignment}
\difficulty{4}
Let $n$ be a natural number, let $g:\{−1,1\}^n\to\{\pm1\}$, and let $O$ be either no
output or a decision tree querying coordinates in $\{1,\ldots,n\}$, with leaves labeled
by signs. If the output, viewed as a tree allowed to query all coordinates, fails to
compute $g$ on the entire Boolean cube whenever it is present, then $O$ is bad for
computing $g$ under the empty fixed-coordinate assignment: when present it fails to
compute $g$ on inputs matching that assignment, and when absent it is deemed bad by
definition.
\end{theorem}

\begin{theorem}[Failure probability of the recursive junta learner]
\lean{FABL.recursiveJuntaLearner_failureProbability_le}
\proofsource{Analysis-of-Boolean-Functions-by-Ryan-ODonnell}{
  pdf-ffbe9c7b3306-p156-b002,
  pdf-ffbe9c7b3306-p156-b004,
  pdf-ffbe9c7b3306-p156-b005,
  pdf-ffbe9c7b3306-p156-b008,
  pdf-ffbe9c7b3306-p168-b001
}
\label{FABL.recursiveJuntaLearner_failureProbability_le}
\leanok
\uses{FABL.BooleanFunction, FABL.F₂DecisionTree.castAvailable, FABL.IsKJunta, FABL.JuntaLearnerOutputBad, FABL.JuntaRelevantCoordinateFinder, FABL.JuntaTreeOutputBad, FABL.LearningProgram.eventProbability, FABL.PositiveLearningParameter, FABL.emptyJuntaFixedAssignment, FABL.eventProbability_map_output_le, FABL.juntaTreeCallCount, FABL.juntaTreeCallCount_mul_perCallFailure, FABL.juntaTreePerCallFailure, FABL.recursiveJuntaLearner, FABL.recursiveJuntaLearnerAux, FABL.recursiveJuntaLearnerAux_failureProbability_le, FABL.rootOutput_bad_implies_auxBad}
\difficulty{4}
Let $n,k\in\mathbb{N}$, let a relevant-coordinate finder be given for dimension $n$ and
junta size $k$, and let $f:\{\pm1\}^n\to\{\pm1\}$ depend on at most $k$ coordinates. For
every $\delta\in(0,1/2]$, the probability that the recursive learner either returns no
tree or returns a decision tree that does not compute $f$ is at most $\delta$.
\end{theorem}

\begin{theorem}[Depth of the recursive junta learner]
\lean{FABL.recursiveJuntaLearner_depth_le}
\label{FABL.recursiveJuntaLearner_depth_le}
\leanok
\uses{FABL.BooleanFunction, FABL.DecisionTree, FABL.F₂DecisionTree.castAvailable, FABL.F₂DecisionTree.depth, FABL.F₂DecisionTree.depth_castAvailable, FABL.JuntaRelevantCoordinateFinder, FABL.LearningCost, FABL.LearningProgram.runWithCost, FABL.LearningProgram.runWithCost_map, FABL.PositiveLearningParameter, FABL.Sign, FABL.emptyJuntaFixedAssignment, FABL.juntaTreePerCallFailure, FABL.recursiveJuntaLearner, FABL.recursiveJuntaLearnerAux_depth_le}
\difficulty{5}
Let $n,k\in\mathbb{N}$, let a relevant-coordinate finder be given for dimension $n$ and
junta size $k$, let $f:\{\pm1\}^n\to\{\pm1\}$ be a Boolean function, and let
$\delta\in(0,1/2]$. If an outcome in the support of the recursive learner's run on $f$
has a decision tree $T$ as its output, then the depth of $T$ is at most $k$.
\end{theorem}

\begin{theorem}[Recursive junta learner cost bounds]
\lean{FABL.recursiveJuntaLearner_cost_le}
\proofsource{Analysis-of-Boolean-Functions-by-Ryan-ODonnell}{
  pdf-ffbe9c7b3306-p156-b005,
  pdf-ffbe9c7b3306-p168-b001
}
\label{FABL.recursiveJuntaLearner_cost_le}
\leanok
\uses{FABL.BooleanFunction, FABL.DecisionTree, FABL.JuntaRelevantCoordinateFinder, FABL.LearningCost, FABL.LearningProgram.runWithCost, FABL.LearningProgram.runWithCost_map, FABL.PositiveLearningParameter, FABL.Sign, FABL.emptyJuntaFixedAssignment, FABL.juntaTreeCallCount, FABL.juntaTreePerCallFailure, FABL.recursiveJuntaLearner, FABL.recursiveJuntaLearnerAux_queries_eq_zero, FABL.recursiveJuntaLearnerAux_randomExamples_le, FABL.recursiveJuntaLearnerAux_work_le}
\difficulty{3}
Let $n,k\in\mathbb{N}$, let a relevant-coordinate finder for dimension $n$ and junta
size $k$ provide a random-example learner with declared sample and work bounds, and let
$f:\{\pm1\}^n\to\{\pm1\}$ be a Boolean function. For a failure parameter
$\delta\in(0,1/2]$, every outcome in the support of the recursive learner's run on $f$
uses at most $T_k B(\delta/T_k)$ random examples, makes no membership queries, and has
charged work at most $T_k W(\delta/T_k)$, where $T_0=1$, $T_{r+1}=1+2T_r$, and $B,W$ are
the finder's declared sample and work bounds.
\end{theorem}
